% readme_fig1.tex: the paper's Fig. 1 for the README, on a transparent page, in a light and a dark version.
% Built by scripts/make_readme_fig1.py, which runs from figures/:
%   light: pdflatex -jobname=readme_fig1_light readme_fig1.tex
%   dark:  pdflatex -jobname=readme_fig1_dark "\def\readmedark{}% readme_fig1.tex: the paper's Fig. 1 for the README, on a transparent page, in a light and a dark version.
% Built by scripts/make_readme_fig1.py, which runs from figures/:
%   light: pdflatex -jobname=readme_fig1_light readme_fig1.tex
%   dark:  pdflatex -jobname=readme_fig1_dark "\def\readmedark{}% readme_fig1.tex: the paper's Fig. 1 for the README, on a transparent page, in a light and a dark version.
% Built by scripts/make_readme_fig1.py, which runs from figures/:
%   light: pdflatex -jobname=readme_fig1_light readme_fig1.tex
%   dark:  pdflatex -jobname=readme_fig1_dark "\def\readmedark{}% readme_fig1.tex: the paper's Fig. 1 for the README, on a transparent page, in a light and a dark version.
% Built by scripts/make_readme_fig1.py, which runs from figures/:
%   light: pdflatex -jobname=readme_fig1_light readme_fig1.tex
%   dark:  pdflatex -jobname=readme_fig1_dark "\def\readmedark{}\input{readme_fig1.tex}"
% The dark colors follow the README's other figures (THEMES["dark"] in scripts/make_readme_assets.py).
\documentclass[conference]{IEEEtran}
\usepackage{amsmath,amssymb}
\usepackage{xcolor}
\usepackage{pgfplots}
\pgfplotsset{compat=1.18}
\usepgfplotslibrary{fillbetween}
\usetikzlibrary{arrows.meta}
\input{icorr_style.tex}
\ifdefined\readmedark
  \definecolor{icInk}{HTML}{E6EDF3}
  \definecolor{icMuted}{HTML}{9DA7B3}
  \definecolor{icRule}{HTML}{30363D}
  \definecolor{icBlue}{HTML}{8FA4FD}
  \definecolor{icBlueLight}{HTML}{3F57C8}
  \definecolor{icBlueDark}{HTML}{C3CEFF}
  \definecolor{icGeneric}{HTML}{8B949E}
  \definecolor{icContext}{HTML}{4B535D}
  \definecolor{icGrey}{HTML}{9DA7B3}
  \definecolor{icDark}{HTML}{C9D1D9}
  \definecolor{icPairA}{HTML}{E58A9A}
  \definecolor{icPairB}{HTML}{33BBEE}
  \definecolor{icOrangeText}{HTML}{F8B860}
  \definecolor{icTintBlue}{HTML}{161D33}% PsyPref's and the user's box
  \definecolor{white}{HTML}{161B22}% the other boxes, slightly raised from GitHub's dark background
\fi
\newcommand{\icorrtikzdata}{data/}
\input{data/fig3p_text.tex}
\input{data/fig4m_text.tex}
\pagestyle{empty}
\begin{document}
\thispagestyle{empty}
\noindent{\color{icInk}\input{fig1.tex}}
\end{document}
"
% The dark colors follow the README's other figures (THEMES["dark"] in scripts/make_readme_assets.py).
\documentclass[conference]{IEEEtran}
\usepackage{amsmath,amssymb}
\usepackage{xcolor}
\usepackage{pgfplots}
\pgfplotsset{compat=1.18}
\usepgfplotslibrary{fillbetween}
\usetikzlibrary{arrows.meta}
% icorr_style.tex: shared colours and line styles of the amplitude-free paper's figures
% (paper/notes/story_concept.md section 7a, scheme C). \input once in the preamble, after
% \usepackage{pgfplots}. Every figure uses these names; no figure defines its own colours.
%
% Roles:
%   PsyPref (all variants) ............ icBlue, solid
%   best published comparator ......... icOrange, solid, hollow marker: square CMA-ES, triangle Arens' GP
%   GP-BALD (same prior) .............. icGrey, dashed
%   CMA-ES on Schaefer's family ....... icDark, dash-dot
%   generic assistance (50 %) ......... icGeneric, solid thin
%   best K-weight profile (ceiling, floor) icContext, dotted
%   K = 8 / 16 / 24 (or k = 24 / 40 / 150) icBlueLight / icBlue / icBlueDark
%   settings as markers ............... circle 3-35 Nm, square fixed 12 Nm, diamond 3-12 Nm
%   robustness: filled = fixed yardstick (post hoc), hollow = registered scoring
\definecolor{icBlue}{HTML}{1845FB}
\definecolor{icBlueLight}{HTML}{8FA4FD}
\definecolor{icBlueDark}{HTML}{0F2C9E}
\definecolor{icOrange}{HTML}{F89C20}
\definecolor{icOrangeText}{HTML}{B86E00}% text in the comparator's hue (contrast 4.0 on white)
\definecolor{icGrey}{HTML}{7A7A7A}
\definecolor{icDark}{HTML}{333333}
\definecolor{icGeneric}{HTML}{9C9CA1}
\definecolor{icContext}{HTML}{BBBBBB}
\definecolor{icRule}{HTML}{DDDDDD}
\definecolor{icInk}{HTML}{222222}
\definecolor{icMuted}{HTML}{666666}
% Fig. 1 only: the pair shown to the user (author's choice, 2026-10-29: distinct hues read
% stronger than two blues; gold was dropped because it equals the comparator's orange)
\definecolor{icTintBlue}{HTML}{F1F3F8}% Fig. 1: the loop, PsyPref's and the user's box (author's choice)
\definecolor{icPairA}{HTML}{BB5566}
\definecolor{icPairB}{HTML}{33BBEE}
\tikzset{
  ic psypref/.style={icBlue, line width=1.5pt, line cap=round, line join=round},
  ic best/.style={icOrange, line width=1.2pt, line cap=round, line join=round},
  ic gpbald/.style={icGrey, line width=0.9pt, dash pattern=on 3.2pt off 1.6pt},
  ic sch/.style={icDark, line width=0.9pt, dash pattern=on 3.6pt off 1.4pt on 0.9pt off 1.4pt},
  ic generic/.style={icGeneric, line width=0.6pt},
  ic context/.style={icContext, line width=0.8pt, dash pattern=on 0.9pt off 1.5pt},
  ic rule/.style={icRule, line width=0.5pt},
  ic label/.style={inner sep=0.6pt, font=\scriptsize, align=left},
  ic note/.style={inner sep=0.8pt, font=\scriptsize, icInk, fill=white, fill opacity=0.85,
                  text opacity=1, align=right},
  ic band/.style={fill=#1, fill opacity=0.15, draw=none},
}
\pgfplotsset{
  ic axis/.style={
    axis line style={icMuted, line width=0.5pt},
    tick style={icMuted, line width=0.5pt},
    tick label style={font=\scriptsize, icMuted},
    label style={font=\footnotesize, icInk},
    title style={font=\footnotesize, icInk},
    axis x line*=bottom, axis y line*=left,
    tick align=outside, major tick length=2pt,
    clip=false,
  },
}

\ifdefined\readmedark
  \definecolor{icInk}{HTML}{E6EDF3}
  \definecolor{icMuted}{HTML}{9DA7B3}
  \definecolor{icRule}{HTML}{30363D}
  \definecolor{icBlue}{HTML}{8FA4FD}
  \definecolor{icBlueLight}{HTML}{3F57C8}
  \definecolor{icBlueDark}{HTML}{C3CEFF}
  \definecolor{icGeneric}{HTML}{8B949E}
  \definecolor{icContext}{HTML}{4B535D}
  \definecolor{icGrey}{HTML}{9DA7B3}
  \definecolor{icDark}{HTML}{C9D1D9}
  \definecolor{icPairA}{HTML}{E58A9A}
  \definecolor{icPairB}{HTML}{33BBEE}
  \definecolor{icOrangeText}{HTML}{F8B860}
  \definecolor{icTintBlue}{HTML}{161D33}% PsyPref's and the user's box
  \definecolor{white}{HTML}{161B22}% the other boxes, slightly raised from GitHub's dark background
\fi
\newcommand{\icorrtikzdata}{data/}
% fig3p_text.tex: written by scripts/paper_v2_fig3p.py (repository); do not edit by hand.
% Every number of Fig. 3's caption and the draft Results paragraph, formatted from values the script checked
% (fam_report.md, fam_numbers.json, amp_table1.json, Fig. 2's data). Differences as fam_report.md prints them.
\def\pjK{24}
\def\pjNusers{200}
\def\pjNsubj{5}
\def\pjSigma{1}
\def\pjLapse{0.1}
\def\pjGapW{18.8}
\def\pjGapWci{(14.7, 22.8)}
\def\pjGapF{9.0}
\def\pjGapFci{(6.1, 11.9)}
\def\pjSchCeilW{94.1}
\def\pjSchCeilF{88.9}
\def\pjSchPsyW{86.4}
\def\pjSchPsyF{74.8}
\def\pjCmaW{67.6}
\def\pjCmaWci{(63.1, 71.9)}
\def\pjCmaF{65.8}
\def\pjCmaFci{(63.5, 67.9)}
\def\pjPubSch{76.6}
\def\pjPubSchci{(74.6, 78.5)}
\def\pjPubPchip{86.6}
\def\pjPubPchipci{(85.0, 88.3)}
\def\pjPubSchDiff{17.5}
\def\pjPubSchDiffci{(15.6, 19.5)}
\def\pjPubPchipDiff{7.4}
\def\pjPubPchipDiffci{(5.9, 9.0)}
\def\pjProjSchW{6.1}
\def\pjProjSchWci{($-7.8$, $-4.6$)}
\def\pjProjSchWsigned{$-6.1$}
\def\pjProjFranksW{$-0.6$}
\def\pjProjPchipW{$-1.1$}
\def\pjProjFixLo{3.4}
\def\pjProjFixHi{9.8}
\def\pjPrompPsyW{92.5}
\def\pjPubTimeFranksW{91.8}
\def\pjPubTimeFranksF{74.2}
\def\pjPubTimePchipW{93.9}
\def\pjPubTimePchipF{89.1}
\def\pjExKpair{15}

% fig4m_text.tex: written by scripts/paper_v2_fig4m.py (repository); do not edit by hand.
\def\fmSThreeW{$+16.4$ ($+14.4$, $+18.3$)}
\def\fmSThreeF{$+7.9$ ($+6.1$, $+9.2$)}
\def\fmSThreeLevW{75.3}
\def\fmSThreeLevF{63.6}
\def\fmStartC{1.5}
\def\fmStartP{2.8}
\def\fmK{24}
\def\fmNclear{68}
\def\fmNall{75}
\def\fmNcell{70}
\def\fmNfix{68}
\def\fmNusers{200}
\def\fmNother{120}
\def\fmNtie{160}
\def\fmMainW{$+13.1$ ($+10.3$, $+14.7$)}
\def\fmMainF{$+8.8$ ($+6.8$, $+10.5$)}
\def\fmSingleW{$+2.6$ ($-0.5$, $+4.8$)}
\def\fmSingleF{$+3.4$ ($+0.5$, $+6.2$)}
\def\fmOtherW{$+16.0$ ($+12.6$, $+19.2$)}
\def\fmOtherF{$+12.0$ ($+8.5$, $+14.3$)}
\def\fmPeaks{$+8.1$}
\def\fmLaterW{$+12.1$}
\def\fmLaterF{$+10.7$}
\def\fmNoneF{$-2.0$ ($-5.3$, $+1.0$)}
\def\fmShapesF{$-3.1$ ($-4.8$, $-1.4$)}
\def\fmShapesW{$+13.7$ ($+10.8$, $+15.7$)}
\def\fmShapesNext{$+15.6$}
\def\fmOtherShapesW{$+5.0$}
\def\fmOtherShapesF{$+5.1$}
\def\fmPessW{$+12.6$ ($+10.0$, $+14.6$)}
\def\fmPessF{$-1.3$ ($-2.9$, $+0.3$)}
\def\fmDropSingle{17.2}
\def\fmDropPess{21.9}
\def\fmFixGap{8.6}
\def\fmBaldLo{1.9}
\def\fmBaldHi{5.7}

\pagestyle{empty}
\begin{document}
\thispagestyle{empty}
\noindent{\color{icInk}% fig1.tex: Fig. 1 of the amplitude-free paper (the method diagram), in TikZ: 0.95 x 252 = 239.4 pt
% wide and 177.84 pt + \WEd in = 186.05 pt tall, for
%   \input{figures/tikz_v2/fig1.tex}
% The simulated users' box shows a worked answer: the pair the loop asks, its distances to the
% preferred profile and the probability they give.
% Preamble: the paper's (pgfplots, compat 1.18, arrows.meta, fillbetween) plus
% \input{icorr_style.tex}; colours are icorr_style.tex's names only.
% Data: \icorrtikzdata (default figures/tikz_v2/data/), fig1_*.csv (scripts/paper_v2_fig1.py,
% repository): one PsyPref-T session in the wide setting (3-35 Nm) for the lower of the two
% middle users of FULL's 200 by distance after 24 comparisons (rank 100 of 200), rerun as FULL's
% run_one builds it and checked against FULL's stored distances; belief particles and the delivered
% profile at k = 0, 6, 24, the preferred profile, generic assistance, 30 draws of the walking-data
% shape prior (unit peaks), the answer model's sigmoid and the walker silhouette. fig1j_*.csv
% (scripts/paper_v2_fig1j.py, the same session): the pair, by rule the one of the first 24
% comparisons with d(A) < d(B) whose d(B) - d(A) lies closest to 1 JND; its distances and
% P(choose A). The CSV columns keep pair1 and pair2 for A and B, and fig1_sigmoid.csv's p is the
% probability of the first profile.
% Layout: the old Fig. 1's coordinates in inches of the script (3.5 x 2.6 in), scaled as the
% paper scales it: 1 in = 0.95 x 72 pt = \mdu; line widths scaled alike (\mds).
\providecommand{\icorrtikzdata}{figures/tikz_v2/data/}%
\makeatletter
\providecommand{\icorrload}[2]{% \icorrload{name}{file.csv}: every cell as \icv{name}{row}{column}%
  \pgfplotstableread[col sep=comma]{\icorrtikzdata #2}\ic@table
  \pgfplotstableforeachcolumn\ic@table\as\ic@col{%
    \pgfplotstableforeachcolumnelement{\ic@col}\of\ic@table\as\ic@cell{%
      \expandafter\global\expandafter\let\csname ic@#1@\pgfplotstablerow @\ic@col\endcsname
        \ic@cell}}}%
\providecommand{\icv}[3]{\csname ic@#1@#2@#3\endcsname}%
% the walker: \icfwalker{x}{y}{height}{last part}: hip at x, ground at y, height, all three
% lengths; last part 5 draws the walker alone, 11 with the hip exoskeleton. Geometry from
% fig1_walker_part<i>.csv (units of the height), style and colour names from
% fig1_walker_parts.csv.
\providecommand{\icfwalker}[4]{%
  \begin{axis}[at={({#1},{#2})}, anchor=origin, x=#3, y=#3, hide axis, clip=false,%
               xmin=-0.3, xmax=0.3, ymin=0, ymax=1]%
    \pgfplotsinvokeforeach{0,...,#4}{\icf@walkerpart{##1}{#3}}%
  \end{axis}}%
\def\icf@none{none}\def\icf@line{line}%
\providecommand{\icf@walkerpart}[2]{%
  \edef\icf@edge{\icv{fwk}{#1}{edge}}\edef\icf@fill{\icv{fwk}{#1}{fill}}%
  \edef\icf@kind{\icv{fwk}{#1}{kind}}%
  \ifx\icf@edge\icf@none\else\colorlet{icfwe#1}{\icf@edge}\fi
  \ifx\icf@fill\icf@none\else\colorlet{icfwf#1}{\icf@fill}\fi
  \ifx\icf@kind\icf@line
    \addplot[draw=icfwe#1, line width=\icv{fwk}{#1}{lw_per_height}*#2,%
             line cap=\icv{fwk}{#1}{cap}, line join=round, mark=none]%
      table[col sep=comma, x=x, y=y] {\icorrtikzdata fig1_walker_part#1.csv};%
  \else\ifx\icf@edge\icf@none
    \addplot[fill=icfwf#1, draw=none, mark=none]%
      table[col sep=comma, x=x, y=y] {\icorrtikzdata fig1_walker_part#1.csv};%
  \else
    \addplot[fill=icfwf#1, draw=icfwe#1, line width=\icv{fwk}{#1}{lw_per_height}*#2, mark=none]%
      table[col sep=comma, x=x, y=y] {\icorrtikzdata fig1_walker_part#1.csv};%
  \fi\fi}%
\makeatother
\begin{tikzpicture}[x=68.4pt, y=68.4pt, font=\scriptsize]
\scriptsize% the type the legend row is measured in
\def\mdu{68.4pt}% 1 in of the script, as printed
\def\mds{0.95}% the paper's scale: line widths
\icorrload{fm}{fig1_meta.csv}\icorrload{fwk}{fig1_walker_parts.csv}\icorrload{fj}{fig1j_meta.csv}%
\pgfplotstableread[col sep=comma]{\icorrtikzdata fig1_curves.csv}\mdcurves
\pgfplotstableread[col sep=comma]{\icorrtikzdata fig1j_pair.csv}\mdpair
\pgfplotstableread[col sep=comma]{\icorrtikzdata fig1_prior.csv}\mdprior
\pgfplotstableread[col sep=comma]{\icorrtikzdata fig1_fan_k0.csv}\mdfanA
\pgfplotstableread[col sep=comma]{\icorrtikzdata fig1_fan_k6.csv}\mdfanB
\pgfplotstableread[col sep=comma]{\icorrtikzdata fig1_fan_k24.csv}\mdfanC
\pgfplotstableread[col sep=comma]{\icorrtikzdata fig1_sigmoid.csv}\mdsigmoid
\pgfplotstablegetcolsof{\mdfanA}\pgfmathtruncatemacro\mdlastdraw{\pgfplotsretval-2}% d0..dN
\pgfplotstablegetcolsof{\mdprior}\pgfmathtruncatemacro\mdlastprior{\pgfplotsretval-3}% s0..sN
\tikzset{
  % a line of text centred on the given y (its baseline 0.257 em below)
  md text/.style={inner sep=0pt, icInk, anchor=base west, yshift=-0.257*7pt*\mds},
  md head/.style={inner sep=0pt, icInk, anchor=base west, yshift=-0.257*8pt*\mds,
                  font=\footnotesize\bfseries},
  % several lines, centred on y
  md block/.style={inner sep=0pt, icInk, anchor=west, align=left,
                   font=\fontsize{7}{7.7}\selectfont},
  md box/.style={draw=icMuted, fill=white, line width=\mds*0.6pt, rounded corners=0.05*\mdu},
  md arrow/.style={draw=#1, line width=\mds*0.8pt,
                   -{Triangle[length=\mds*2.8pt, width=\mds*2.8pt]}},
  md arrow/.default=icMuted,
  md generic/.style={icGeneric, line width=\mds*1.0pt},
}
\pgfplotsset{
  md mini/.style={anchor=south west, scale only axis, hide axis, clip=false,
    xmin=\icv{fm}{0}{xmin}, xmax=\icv{fm}{0}{xmax},
    ymin=\icv{fm}{0}{ymin}, ymax=\icv{fm}{0}{ymax}},
  md line/.style={line cap=round, line join=round, mark=none},
}
% a small bare axes: gait cycle on x, torque on y, a zero line; #1, #2 lower-left corner,
% #3, #4 width and height (in), #5 its curves; the style md extra: further axis options
\pgfplotsset{md extra/.style={}}
\def\mdmini#1#2#3#4#5{%
  \begin{axis}[md mini, at={({(#1)*\mdu},{(#2)*\mdu})}, width=(#3)*\mdu, height=(#4)*\mdu,
               md extra]
    \draw[icRule, line width=\mds*0.5pt]
      (axis cs:\icv{fm}{0}{xmin},0) -- (axis cs:\icv{fm}{0}{xmax},0);
    #5
    \draw[icContext, line width=\mds*0.6pt]
      (rel axis cs:0,1) -- (rel axis cs:0,0) -- (rel axis cs:1,0);
  \end{axis}}
% a fan of posterior particles (#1 the table, #2 colour, #3 opacity)
\def\mdfan#1#2#3{%
  \pgfplotsinvokeforeach{0,...,\mdlastdraw}{%
    \addplot[md line, #2, line width=\mds*0.4pt, draw opacity=#3] table[x=phase, y=d##1] {#1};}}
% the prior's shape draws
\def\mdpriorfan{%
  \pgfplotsinvokeforeach{0,...,\mdlastprior}{%
    \addplot[md line, icContext, line width=\mds*0.35pt, draw opacity=0.7]
      table[x=phase, y=s##1] {\mdprior};}}
\def\mdpreferred#1{\addplot[md line, icInk, line width=\mds*#1pt] table[x=phase, y=preferred]
  {\mdcurves};}
\def\mdestimate#1#2{\addplot[md line, icBlue, line width=\mds*#2pt] table[x=phase, y=#1]
  {\mdcurves};}
\def\mdgeneric{\addplot[md generic, md line] table[x=phase, y=generic] {\mdcurves};}
% a profile of the pair (#1 its column, #2 colour), the preferred profile, the gap between them
\def\mdgap#1#2{%
  \addplot[name path=mdgapA, draw=none, mark=none] table[x=phase, y=#1] {\mdpair};
  \addplot[name path=mdgapB, draw=none, mark=none] table[x=phase, y=preferred] {\mdpair};
  \addplot[#2, fill opacity=0.25, draw=none] fill between[of=mdgapA and mdgapB];
  \addplot[md line, icInk, line width=\mds*0.8pt] table[x=phase, y=preferred] {\mdpair};
  \addplot[md line, #2, line width=\mds*1.1pt] table[x=phase, y=#1] {\mdpair};}
% columns and rows (in)
\def\LXa{0.03}\def\LXb{0.77}\def\MXa{0.86}\def\MXb{2.58}\def\RXa{2.67}\def\RXb{3.47}
\def\TYa{0.80}\def\TYb{2.46}
\pgfmathsetmacro\lc{(\LXa+\LXb)/2}\pgfmathsetmacro\rc{(\RXa+\RXb)/2}
% ---- legend: one centred row (0.18 in lines, 0.05 in pad, 0.16 in gaps), widths of the labels
% measured in this type
\pgfmathsetmacro\lwa{width("PsyPref")/68.4}
\pgfmathsetmacro\lwb{width("preferred profile")/68.4}
\pgfmathsetmacro\lwc{width("generic assistance")/68.4}
\pgfmathsetmacro\lxa{(3.5-(3*(0.18+0.05)+\lwa+\lwb+\lwc+2*0.16))/2}
\pgfmathsetmacro\lxb{\lxa+0.18+0.05+\lwa+0.16}
\pgfmathsetmacro\lxc{\lxb+0.18+0.05+\lwb+0.16}
\def\LGy{2.548}
\draw[icBlue, line width=\mds*1.2pt] (\lxa,\LGy) -- ++(0.18,0);
\node[md text] at (\lxa+0.23,\LGy) {PsyPref};
\draw[icInk, line width=\mds*0.9pt] (\lxb,\LGy) -- ++(0.18,0);
\node[md text] at (\lxb+0.23,\LGy) {Preferred profile};
\draw[md generic] (\lxc,\LGy) -- ++(0.18,0);
\node[md text] at (\lxc+0.23,\LGy) {Generic assistance};
% ---- left: what the prior is built from <- the study's peak-torque range; italics: where prior
% knowledge enters
\draw[md box] (\LXa,2.07) rectangle (\LXb,\TYb);
\node[md block, anchor=center, align=center] at (\lc,2.19) {Walking data\\\textit{or other curves}};
\draw[md arrow] (\lc,2.065) -- (\lc,1.985);
\draw[md box] (\LXa,1.25) rectangle (\LXb,1.98);
\node[md block, anchor=center, align=center] at (\lc,1.835) {Profile-shape\\prior, \textit{width}};
\node[md text, anchor=base] at (\LXa+0.15+0.27,1.315) {Gait cycle};
\node[inner sep=0pt, icInk, rotate=90, anchor=center] at (\LXa+0.075,1.40+0.14) {Torque};
\draw[icMuted, line width=\mds*0.6pt, {Triangle[length=\mds*2.2pt, width=\mds*2.2pt]}-{Triangle[length=\mds*2.2pt, width=\mds*2.2pt]}]
  (\LXb-0.033,1.42) -- (\LXb-0.033,1.68);
\draw[md arrow] (\lc,1.195) -- (\lc,1.245);
\draw[md box] (\LXa,\TYa) rectangle (\LXb,1.19);
\node[md block, anchor=center, align=center] at (\lc,0.995)
  {Peak torque\\range or level\\\textit{device, protocol}};
\draw[md arrow] (\LXb+0.005,1.89) -- (\MXa-0.005,1.89);
% ---- middle top: PsyPref's belief
\draw[md box, fill=icTintBlue, draw=icBlue, line width=\mds*0.9pt] (\MXa,1.70) rectangle (\MXb,\TYb);
\node[md head, text=icBlue] at (\MXa+0.07,2.375) {PsyPref:};
\node[md text] at (\MXa+0.07,2.258) {belief over the preferred profile $\mathbf z^\ast$,};
\node[md text] at (\MXa+0.07,2.146) {discrimination scale $\ell$, lapse rate $\lambda$};
\pgfmathsetmacro\pxa{(\MXa+\MXb)/2-(3*0.46+2*0.115)/2}
\pgfmathsetmacro\pxb{\pxa+0.46+0.115}\pgfmathsetmacro\pxc{\pxb+0.46+0.115}
\def\SNy{1.835}\def\SNh{0.24}% the belief insets: bottom and height
\node[md text, anchor=base] at (\pxa+0.23,1.77) {$k = \icv{fm}{0}{k0}$};
\node[md text, anchor=base] at (\pxb+0.23,1.77) {$k = \icv{fm}{0}{k1}$};
\node[md text, anchor=base] at (\pxc+0.23,1.77) {$k = \icv{fm}{0}{k2}$};
\draw[md arrow=icContext, line width=\mds*0.7pt]
  (\pxb-0.115+0.025,\SNy+\SNh/2) -- (\pxb-0.025,\SNy+\SNh/2);
\draw[md arrow=icContext, line width=\mds*0.7pt]
  (\pxc-0.115+0.025,\SNy+\SNh/2) -- (\pxc-0.025,\SNy+\SNh/2);
% ---- middle bottom: the user
\draw[md box, fill=icTintBlue] (\MXa,\TYa) rectangle (\MXb,1.40);
\pgfmathsetmacro\qx{\MXa+0.06+0.44+0.08}
\node[md block, font=\fontsize{7}{8}\selectfont\bfseries] at (\qx,1.27)
  {Profile A, B or\\no preference?};
% ---- the loop's two edges
\pgfmathsetmacro\xl{\MXa+0.10}\pgfmathsetmacro\xr{\MXb-0.10}
\draw[md arrow=icInk] (\xr,1.695) -- (\xr,1.405);
\node[md block, anchor=east, align=right] at (\xr-0.05,1.55) {Most informative\\pair};
\draw[md arrow=icInk] (\xl,1.405) -- (\xl,1.695);
\node[md block] at (\xl+0.05,1.55) {Belief\\update};
% ---- right: identified assistance -> validation
\draw[md arrow] (\MXb+0.005,2.01) -- (\RXa-0.005,2.01);
\draw[md box] (\RXa,1.70) rectangle (\RXb,\TYb);
\node[md block] at (\RXa+0.06,2.2024)
  {\textcolor{icBlue}{Identified}\\\textcolor{icBlue}{assistance}\\after $k$\\comparisons};
\draw[md arrow] (\rc,1.695) -- (\rc,1.615);
\draw[md box] (\RXa,\TYa) rectangle (\RXb,1.61);
\node[md text, font=\scriptsize\bfseries] at (\RXa+0.06,1.515) {Validation:};
\node[md block] at (\RXa+0.06,1.30) {identified vs.\\generic\\assistance};
% ---- bottom: this paper's simulated users
% the worked answer: the plots' bottom and height, the pair's left edges and width, the answer
% model's left edge and width, the row of labels below them, how far the box grows downward (in)
\def\WEy{0.07}\def\WEh{0.34}\def\WEa{1.07}\def\WEb{1.68}\def\WEw{0.52}
\def\SGx{2.38}\def\SGw{0.54}\def\SGy{\WEy}\def\SGh{\WEh}
\def\WEl{-0.025}\def\WEd{0.12}
\draw[draw=icMuted, line width=\mds*0.8pt, dash pattern=on \mds*2.4pt off \mds*1.6pt,
      rounded corners=0.06*\mdu] (\LXa,0.03-\WEd) rectangle (\RXb,0.70);
\node[md head] at (\LXa+0.07,0.615) {This paper: simulated users};
% who the users are; how they answer: the pair's distances, then the probability they give
\node[md block, anchor=north west] at (\LXa+0.07,0.535)
  {\textbf{Preferred profile}\\\textit{Shape:} walker not\\\phantom{\textit{Shape:} }in the prior\\\textit{Peaks:} study's range};
\node[md block, anchor=north west] at (\WEa,0.535) {\textbf{Distance} $d$ (Nm)};
\node[md block, anchor=north west] at (2.12,0.535) {\textbf{Answer model:} $P$(choose A)};
\node[md text, anchor=base] at (\WEa+\WEw/2,\WEl) {$d(\mathrm{A})\!=\!\icv{fj}{0}{d1_txt}$};
\node[md text, anchor=base] at (\WEb+\WEw/2,\WEl) {$d(\mathrm{B})\!=\!\icv{fj}{0}{d2_txt}$};
\node[md text, anchor=base] at (\SGx+\SGw/2,\WEl) {$d(\mathrm{B}) - d(\mathrm{A})$};
\draw[md arrow] (\WEb+\WEw+0.03,\WEy+\WEh/2) -- (\SGx-0.05,\WEy+\WEh/2);
\draw[md arrow] (\MXa+0.28,0.705) -- (\MXa+0.28,0.795);
% ---- the walkers: three without, the user with the hip exoskeleton
\icfwalker{(\LXa+0.20)*\mdu}{2.315*\mdu}{0.12*\mdu}{5}
\icfwalker{(\LXa+0.38)*\mdu}{2.315*\mdu}{0.12*\mdu}{5}
\icfwalker{(\LXa+0.56)*\mdu}{2.315*\mdu}{0.12*\mdu}{5}
\icfwalker{(\MXa+0.06+0.22-0.02)*\mdu}{(\TYa+0.05+0.02)*\mdu}{0.46*\mdu}{11}
% ---- the small axes
% the prior's shapes at unit peaks: limits of their own
\pgfplotsset{md extra/.style={ymin=\icv{fm}{0}{prior_ymin}, ymax=\icv{fm}{0}{prior_ymax}}}
\mdmini{\LXa+0.15}{1.40}{0.54}{0.29}{%
  \mdpriorfan\addplot[md generic, md line] table[x=phase, y=generic_shape] {\mdprior};}
\pgfplotsset{md extra/.style={}}
\mdmini{\pxa}{\SNy}{0.46}{\SNh}{%
  \mdfan{\mdfanA}{icBlueLight}{0.55}\mdestimate{mean_k0}{1.0}\mdpreferred{0.8}}
\mdmini{\pxb}{\SNy}{0.46}{\SNh}{%
  \mdfan{\mdfanB}{icBlueLight}{0.55}\mdestimate{mean_k6}{1.0}\mdpreferred{0.8}}
\mdmini{\pxc}{\SNy}{0.46}{\SNh}{%
  \mdfan{\mdfanC}{icBlueLight}{0.55}\mdestimate{mean_k24}{1.0}\mdpreferred{0.8}}
% the pair: each label inside its plot's top-left corner, clear of the other plot
\mdmini{\qx}{0.87}{0.46}{0.27}{%
  \addplot[md line, icPairA, line width=\mds*1.2pt] table[x=phase, y=pair1] {\mdpair};
  \node[inner sep=0pt, anchor=north west, font=\scriptsize\bfseries, icInk, xshift=\mds*2.5pt,
        yshift=\mds*1.5pt] at (rel axis cs:0,1) {A};}
\mdmini{\qx+0.60}{0.87}{0.46}{0.27}{%
  \addplot[md line, icPairB, line width=\mds*1.2pt] table[x=phase, y=pair2] {\mdpair};
  \node[inner sep=0pt, anchor=north west, font=\scriptsize\bfseries, icInk, xshift=\mds*2.5pt,
        yshift=\mds*1.5pt] at (rel axis cs:0,1) {B};}
% the same pair with the preferred profile: limits of its own, so that the gaps fill the plots
\pgfplotsset{md extra/.style={ymin=\icv{fj}{0}{ymin}, ymax=\icv{fj}{0}{ymax}}}
\mdmini{\WEa}{\WEy}{\WEw}{\WEh}{%
  \mdgap{pair1}{icPairA}
  \node[inner sep=0pt, anchor=north west, font=\scriptsize\bfseries, icInk, xshift=\mds*2.5pt,
        yshift=\mds*1.5pt] at (rel axis cs:0,1) {A};}
\mdmini{\WEb}{\WEy}{\WEw}{\WEh}{%
  \mdgap{pair2}{icPairB}
  \node[inner sep=0pt, anchor=north west, font=\scriptsize\bfseries, icInk, xshift=\mds*2.5pt,
        yshift=\mds*1.5pt] at (rel axis cs:0,1) {B};}
\pgfplotsset{md extra/.style={}}
\mdmini{\RXa+0.06}{1.75}{\RXb-\RXa-0.12}{0.21}{%
  \mdpreferred{0.9}\mdestimate{mean_k24}{1.2}
  \node[inner sep=0pt, anchor=south east, icMuted]
    at (rel axis cs:0.99,0.02) {$k = \icv{fm}{0}{k_id}$};}
\mdmini{\RXa+0.06}{\TYa+0.05}{\RXb-\RXa-0.12}{0.27}{%
  \mdgeneric\mdestimate{mean_k24}{1.2}}
% ---- the answer model: lapse floor and ceiling, ticks on the right
\pgfmathsetmacro\mdlo{\icv{fm}{0}{lapse}/2}\pgfmathsetmacro\mdhi{1-\icv{fm}{0}{lapse}/2}
\begin{axis}[at={(\SGx*\mdu,\SGy*\mdu)}, anchor=south west, scale only axis,
  width=\SGw*\mdu, height=\SGh*\mdu, clip=false,
  xmin=\icv{fm}{0}{sig_xmin}, xmax=\icv{fm}{0}{sig_xmax},
  ymin=\icv{fm}{0}{sig_ymin}, ymax=\icv{fm}{0}{sig_ymax},
  axis x line*=bottom, axis y line*=right, axis line style={-, icMuted, line width=\mds*0.6pt},
  tick align=outside, major tick length=\mds*2pt, tick style={icMuted, line width=\mds*0.6pt},
  xtick={0}, xticklabels={}, ytick={\mdlo,0.5,\mdhi},
  yticklabels={$\lambda$/2, 0.5, 1\,$-$\,$\lambda$/2}, yticklabel pos=right,
  yticklabel style={icInk, font=\scriptsize, inner sep=0pt, xshift=\mds*1.5pt},
]
\draw[icContext, line width=\mds*0.6pt, dash pattern=on \mds*0.6pt off \mds*0.9pt]
  (axis cs:\icv{fm}{0}{sig_xmin},\mdlo) -- (axis cs:\icv{fm}{0}{sig_xmax},\mdlo);
\draw[icContext, line width=\mds*0.6pt, dash pattern=on \mds*0.6pt off \mds*0.9pt]
  (axis cs:\icv{fm}{0}{sig_xmin},\mdhi) -- (axis cs:\icv{fm}{0}{sig_xmax},\mdhi);
\draw[icRule, line width=\mds*0.5pt]
  (axis cs:0,\icv{fm}{0}{sig_ymin}) -- (axis cs:0,\icv{fm}{0}{sig_ymax});
\addplot[icInk, line width=\mds*1.2pt, mark=none, line join=round] table[x=delta, y=p]
  {\mdsigmoid};
% the pair's reading: its distance difference, the probability it gives
\draw[icMuted, line width=\mds*0.6pt, dash pattern=on \mds*0.9pt off \mds*0.9pt]
  (axis cs:\icv{fj}{0}{delta},\icv{fm}{0}{sig_ymin}) -- (axis cs:\icv{fj}{0}{delta},\icv{fj}{0}{p})
  -- (axis cs:\icv{fm}{0}{sig_xmin},\icv{fj}{0}{p});
\fill[icInk] (axis cs:\icv{fj}{0}{delta},\icv{fj}{0}{p}) circle[radius=\mds*1.3pt];
\node[inner sep=0pt, anchor=north west, font=\scriptsize\bfseries, icInk, yshift=-\mds*1.2pt]
  at (axis cs:\icv{fm}{0}{sig_xmin},\icv{fj}{0}{p}) {\icv{fj}{0}{p_txt}};
\end{axis}
\pgfresetboundingbox
\useasboundingbox (0,-\WEd) rectangle (3.5,2.6);
\end{tikzpicture}%
}
\end{document}
"
% The dark colors follow the README's other figures (THEMES["dark"] in scripts/make_readme_assets.py).
\documentclass[conference]{IEEEtran}
\usepackage{amsmath,amssymb}
\usepackage{xcolor}
\usepackage{pgfplots}
\pgfplotsset{compat=1.18}
\usepgfplotslibrary{fillbetween}
\usetikzlibrary{arrows.meta}
% icorr_style.tex: shared colours and line styles of the amplitude-free paper's figures
% (paper/notes/story_concept.md section 7a, scheme C). \input once in the preamble, after
% \usepackage{pgfplots}. Every figure uses these names; no figure defines its own colours.
%
% Roles:
%   PsyPref (all variants) ............ icBlue, solid
%   best published comparator ......... icOrange, solid, hollow marker: square CMA-ES, triangle Arens' GP
%   GP-BALD (same prior) .............. icGrey, dashed
%   CMA-ES on Schaefer's family ....... icDark, dash-dot
%   generic assistance (50 %) ......... icGeneric, solid thin
%   best K-weight profile (ceiling, floor) icContext, dotted
%   K = 8 / 16 / 24 (or k = 24 / 40 / 150) icBlueLight / icBlue / icBlueDark
%   settings as markers ............... circle 3-35 Nm, square fixed 12 Nm, diamond 3-12 Nm
%   robustness: filled = fixed yardstick (post hoc), hollow = registered scoring
\definecolor{icBlue}{HTML}{1845FB}
\definecolor{icBlueLight}{HTML}{8FA4FD}
\definecolor{icBlueDark}{HTML}{0F2C9E}
\definecolor{icOrange}{HTML}{F89C20}
\definecolor{icOrangeText}{HTML}{B86E00}% text in the comparator's hue (contrast 4.0 on white)
\definecolor{icGrey}{HTML}{7A7A7A}
\definecolor{icDark}{HTML}{333333}
\definecolor{icGeneric}{HTML}{9C9CA1}
\definecolor{icContext}{HTML}{BBBBBB}
\definecolor{icRule}{HTML}{DDDDDD}
\definecolor{icInk}{HTML}{222222}
\definecolor{icMuted}{HTML}{666666}
% Fig. 1 only: the pair shown to the user (author's choice, 2026-10-29: distinct hues read
% stronger than two blues; gold was dropped because it equals the comparator's orange)
\definecolor{icTintBlue}{HTML}{F1F3F8}% Fig. 1: the loop, PsyPref's and the user's box (author's choice)
\definecolor{icPairA}{HTML}{BB5566}
\definecolor{icPairB}{HTML}{33BBEE}
\tikzset{
  ic psypref/.style={icBlue, line width=1.5pt, line cap=round, line join=round},
  ic best/.style={icOrange, line width=1.2pt, line cap=round, line join=round},
  ic gpbald/.style={icGrey, line width=0.9pt, dash pattern=on 3.2pt off 1.6pt},
  ic sch/.style={icDark, line width=0.9pt, dash pattern=on 3.6pt off 1.4pt on 0.9pt off 1.4pt},
  ic generic/.style={icGeneric, line width=0.6pt},
  ic context/.style={icContext, line width=0.8pt, dash pattern=on 0.9pt off 1.5pt},
  ic rule/.style={icRule, line width=0.5pt},
  ic label/.style={inner sep=0.6pt, font=\scriptsize, align=left},
  ic note/.style={inner sep=0.8pt, font=\scriptsize, icInk, fill=white, fill opacity=0.85,
                  text opacity=1, align=right},
  ic band/.style={fill=#1, fill opacity=0.15, draw=none},
}
\pgfplotsset{
  ic axis/.style={
    axis line style={icMuted, line width=0.5pt},
    tick style={icMuted, line width=0.5pt},
    tick label style={font=\scriptsize, icMuted},
    label style={font=\footnotesize, icInk},
    title style={font=\footnotesize, icInk},
    axis x line*=bottom, axis y line*=left,
    tick align=outside, major tick length=2pt,
    clip=false,
  },
}

\ifdefined\readmedark
  \definecolor{icInk}{HTML}{E6EDF3}
  \definecolor{icMuted}{HTML}{9DA7B3}
  \definecolor{icRule}{HTML}{30363D}
  \definecolor{icBlue}{HTML}{8FA4FD}
  \definecolor{icBlueLight}{HTML}{3F57C8}
  \definecolor{icBlueDark}{HTML}{C3CEFF}
  \definecolor{icGeneric}{HTML}{8B949E}
  \definecolor{icContext}{HTML}{4B535D}
  \definecolor{icGrey}{HTML}{9DA7B3}
  \definecolor{icDark}{HTML}{C9D1D9}
  \definecolor{icPairA}{HTML}{E58A9A}
  \definecolor{icPairB}{HTML}{33BBEE}
  \definecolor{icOrangeText}{HTML}{F8B860}
  \definecolor{icTintBlue}{HTML}{161D33}% PsyPref's and the user's box
  \definecolor{white}{HTML}{161B22}% the other boxes, slightly raised from GitHub's dark background
\fi
\newcommand{\icorrtikzdata}{data/}
% fig3p_text.tex: written by scripts/paper_v2_fig3p.py (repository); do not edit by hand.
% Every number of Fig. 3's caption and the draft Results paragraph, formatted from values the script checked
% (fam_report.md, fam_numbers.json, amp_table1.json, Fig. 2's data). Differences as fam_report.md prints them.
\def\pjK{24}
\def\pjNusers{200}
\def\pjNsubj{5}
\def\pjSigma{1}
\def\pjLapse{0.1}
\def\pjGapW{18.8}
\def\pjGapWci{(14.7, 22.8)}
\def\pjGapF{9.0}
\def\pjGapFci{(6.1, 11.9)}
\def\pjSchCeilW{94.1}
\def\pjSchCeilF{88.9}
\def\pjSchPsyW{86.4}
\def\pjSchPsyF{74.8}
\def\pjCmaW{67.6}
\def\pjCmaWci{(63.1, 71.9)}
\def\pjCmaF{65.8}
\def\pjCmaFci{(63.5, 67.9)}
\def\pjPubSch{76.6}
\def\pjPubSchci{(74.6, 78.5)}
\def\pjPubPchip{86.6}
\def\pjPubPchipci{(85.0, 88.3)}
\def\pjPubSchDiff{17.5}
\def\pjPubSchDiffci{(15.6, 19.5)}
\def\pjPubPchipDiff{7.4}
\def\pjPubPchipDiffci{(5.9, 9.0)}
\def\pjProjSchW{6.1}
\def\pjProjSchWci{($-7.8$, $-4.6$)}
\def\pjProjSchWsigned{$-6.1$}
\def\pjProjFranksW{$-0.6$}
\def\pjProjPchipW{$-1.1$}
\def\pjProjFixLo{3.4}
\def\pjProjFixHi{9.8}
\def\pjPrompPsyW{92.5}
\def\pjPubTimeFranksW{91.8}
\def\pjPubTimeFranksF{74.2}
\def\pjPubTimePchipW{93.9}
\def\pjPubTimePchipF{89.1}
\def\pjExKpair{15}

% fig4m_text.tex: written by scripts/paper_v2_fig4m.py (repository); do not edit by hand.
\def\fmSThreeW{$+16.4$ ($+14.4$, $+18.3$)}
\def\fmSThreeF{$+7.9$ ($+6.1$, $+9.2$)}
\def\fmSThreeLevW{75.3}
\def\fmSThreeLevF{63.6}
\def\fmStartC{1.5}
\def\fmStartP{2.8}
\def\fmK{24}
\def\fmNclear{68}
\def\fmNall{75}
\def\fmNcell{70}
\def\fmNfix{68}
\def\fmNusers{200}
\def\fmNother{120}
\def\fmNtie{160}
\def\fmMainW{$+13.1$ ($+10.3$, $+14.7$)}
\def\fmMainF{$+8.8$ ($+6.8$, $+10.5$)}
\def\fmSingleW{$+2.6$ ($-0.5$, $+4.8$)}
\def\fmSingleF{$+3.4$ ($+0.5$, $+6.2$)}
\def\fmOtherW{$+16.0$ ($+12.6$, $+19.2$)}
\def\fmOtherF{$+12.0$ ($+8.5$, $+14.3$)}
\def\fmPeaks{$+8.1$}
\def\fmLaterW{$+12.1$}
\def\fmLaterF{$+10.7$}
\def\fmNoneF{$-2.0$ ($-5.3$, $+1.0$)}
\def\fmShapesF{$-3.1$ ($-4.8$, $-1.4$)}
\def\fmShapesW{$+13.7$ ($+10.8$, $+15.7$)}
\def\fmShapesNext{$+15.6$}
\def\fmOtherShapesW{$+5.0$}
\def\fmOtherShapesF{$+5.1$}
\def\fmPessW{$+12.6$ ($+10.0$, $+14.6$)}
\def\fmPessF{$-1.3$ ($-2.9$, $+0.3$)}
\def\fmDropSingle{17.2}
\def\fmDropPess{21.9}
\def\fmFixGap{8.6}
\def\fmBaldLo{1.9}
\def\fmBaldHi{5.7}

\pagestyle{empty}
\begin{document}
\thispagestyle{empty}
\noindent{\color{icInk}% fig1.tex: Fig. 1 of the amplitude-free paper (the method diagram), in TikZ: 0.95 x 252 = 239.4 pt
% wide and 177.84 pt + \WEd in = 186.05 pt tall, for
%   \input{figures/tikz_v2/fig1.tex}
% The simulated users' box shows a worked answer: the pair the loop asks, its distances to the
% preferred profile and the probability they give.
% Preamble: the paper's (pgfplots, compat 1.18, arrows.meta, fillbetween) plus
% % icorr_style.tex: shared colours and line styles of the amplitude-free paper's figures
% (paper/notes/story_concept.md section 7a, scheme C). \input once in the preamble, after
% \usepackage{pgfplots}. Every figure uses these names; no figure defines its own colours.
%
% Roles:
%   PsyPref (all variants) ............ icBlue, solid
%   best published comparator ......... icOrange, solid, hollow marker: square CMA-ES, triangle Arens' GP
%   GP-BALD (same prior) .............. icGrey, dashed
%   CMA-ES on Schaefer's family ....... icDark, dash-dot
%   generic assistance (50 %) ......... icGeneric, solid thin
%   best K-weight profile (ceiling, floor) icContext, dotted
%   K = 8 / 16 / 24 (or k = 24 / 40 / 150) icBlueLight / icBlue / icBlueDark
%   settings as markers ............... circle 3-35 Nm, square fixed 12 Nm, diamond 3-12 Nm
%   robustness: filled = fixed yardstick (post hoc), hollow = registered scoring
\definecolor{icBlue}{HTML}{1845FB}
\definecolor{icBlueLight}{HTML}{8FA4FD}
\definecolor{icBlueDark}{HTML}{0F2C9E}
\definecolor{icOrange}{HTML}{F89C20}
\definecolor{icOrangeText}{HTML}{B86E00}% text in the comparator's hue (contrast 4.0 on white)
\definecolor{icGrey}{HTML}{7A7A7A}
\definecolor{icDark}{HTML}{333333}
\definecolor{icGeneric}{HTML}{9C9CA1}
\definecolor{icContext}{HTML}{BBBBBB}
\definecolor{icRule}{HTML}{DDDDDD}
\definecolor{icInk}{HTML}{222222}
\definecolor{icMuted}{HTML}{666666}
% Fig. 1 only: the pair shown to the user (author's choice, 2026-10-29: distinct hues read
% stronger than two blues; gold was dropped because it equals the comparator's orange)
\definecolor{icTintBlue}{HTML}{F1F3F8}% Fig. 1: the loop, PsyPref's and the user's box (author's choice)
\definecolor{icPairA}{HTML}{BB5566}
\definecolor{icPairB}{HTML}{33BBEE}
\tikzset{
  ic psypref/.style={icBlue, line width=1.5pt, line cap=round, line join=round},
  ic best/.style={icOrange, line width=1.2pt, line cap=round, line join=round},
  ic gpbald/.style={icGrey, line width=0.9pt, dash pattern=on 3.2pt off 1.6pt},
  ic sch/.style={icDark, line width=0.9pt, dash pattern=on 3.6pt off 1.4pt on 0.9pt off 1.4pt},
  ic generic/.style={icGeneric, line width=0.6pt},
  ic context/.style={icContext, line width=0.8pt, dash pattern=on 0.9pt off 1.5pt},
  ic rule/.style={icRule, line width=0.5pt},
  ic label/.style={inner sep=0.6pt, font=\scriptsize, align=left},
  ic note/.style={inner sep=0.8pt, font=\scriptsize, icInk, fill=white, fill opacity=0.85,
                  text opacity=1, align=right},
  ic band/.style={fill=#1, fill opacity=0.15, draw=none},
}
\pgfplotsset{
  ic axis/.style={
    axis line style={icMuted, line width=0.5pt},
    tick style={icMuted, line width=0.5pt},
    tick label style={font=\scriptsize, icMuted},
    label style={font=\footnotesize, icInk},
    title style={font=\footnotesize, icInk},
    axis x line*=bottom, axis y line*=left,
    tick align=outside, major tick length=2pt,
    clip=false,
  },
}
; colours are icorr_style.tex's names only.
% Data: \icorrtikzdata (default figures/tikz_v2/data/), fig1_*.csv (scripts/paper_v2_fig1.py,
% repository): one PsyPref-T session in the wide setting (3-35 Nm) for the lower of the two
% middle users of FULL's 200 by distance after 24 comparisons (rank 100 of 200), rerun as FULL's
% run_one builds it and checked against FULL's stored distances; belief particles and the delivered
% profile at k = 0, 6, 24, the preferred profile, generic assistance, 30 draws of the walking-data
% shape prior (unit peaks), the answer model's sigmoid and the walker silhouette. fig1j_*.csv
% (scripts/paper_v2_fig1j.py, the same session): the pair, by rule the one of the first 24
% comparisons with d(A) < d(B) whose d(B) - d(A) lies closest to 1 JND; its distances and
% P(choose A). The CSV columns keep pair1 and pair2 for A and B, and fig1_sigmoid.csv's p is the
% probability of the first profile.
% Layout: the old Fig. 1's coordinates in inches of the script (3.5 x 2.6 in), scaled as the
% paper scales it: 1 in = 0.95 x 72 pt = \mdu; line widths scaled alike (\mds).
\providecommand{\icorrtikzdata}{figures/tikz_v2/data/}%
\makeatletter
\providecommand{\icorrload}[2]{% \icorrload{name}{file.csv}: every cell as \icv{name}{row}{column}%
  \pgfplotstableread[col sep=comma]{\icorrtikzdata #2}\ic@table
  \pgfplotstableforeachcolumn\ic@table\as\ic@col{%
    \pgfplotstableforeachcolumnelement{\ic@col}\of\ic@table\as\ic@cell{%
      \expandafter\global\expandafter\let\csname ic@#1@\pgfplotstablerow @\ic@col\endcsname
        \ic@cell}}}%
\providecommand{\icv}[3]{\csname ic@#1@#2@#3\endcsname}%
% the walker: \icfwalker{x}{y}{height}{last part}: hip at x, ground at y, height, all three
% lengths; last part 5 draws the walker alone, 11 with the hip exoskeleton. Geometry from
% fig1_walker_part<i>.csv (units of the height), style and colour names from
% fig1_walker_parts.csv.
\providecommand{\icfwalker}[4]{%
  \begin{axis}[at={({#1},{#2})}, anchor=origin, x=#3, y=#3, hide axis, clip=false,%
               xmin=-0.3, xmax=0.3, ymin=0, ymax=1]%
    \pgfplotsinvokeforeach{0,...,#4}{\icf@walkerpart{##1}{#3}}%
  \end{axis}}%
\def\icf@none{none}\def\icf@line{line}%
\providecommand{\icf@walkerpart}[2]{%
  \edef\icf@edge{\icv{fwk}{#1}{edge}}\edef\icf@fill{\icv{fwk}{#1}{fill}}%
  \edef\icf@kind{\icv{fwk}{#1}{kind}}%
  \ifx\icf@edge\icf@none\else\colorlet{icfwe#1}{\icf@edge}\fi
  \ifx\icf@fill\icf@none\else\colorlet{icfwf#1}{\icf@fill}\fi
  \ifx\icf@kind\icf@line
    \addplot[draw=icfwe#1, line width=\icv{fwk}{#1}{lw_per_height}*#2,%
             line cap=\icv{fwk}{#1}{cap}, line join=round, mark=none]%
      table[col sep=comma, x=x, y=y] {\icorrtikzdata fig1_walker_part#1.csv};%
  \else\ifx\icf@edge\icf@none
    \addplot[fill=icfwf#1, draw=none, mark=none]%
      table[col sep=comma, x=x, y=y] {\icorrtikzdata fig1_walker_part#1.csv};%
  \else
    \addplot[fill=icfwf#1, draw=icfwe#1, line width=\icv{fwk}{#1}{lw_per_height}*#2, mark=none]%
      table[col sep=comma, x=x, y=y] {\icorrtikzdata fig1_walker_part#1.csv};%
  \fi\fi}%
\makeatother
\begin{tikzpicture}[x=68.4pt, y=68.4pt, font=\scriptsize]
\scriptsize% the type the legend row is measured in
\def\mdu{68.4pt}% 1 in of the script, as printed
\def\mds{0.95}% the paper's scale: line widths
\icorrload{fm}{fig1_meta.csv}\icorrload{fwk}{fig1_walker_parts.csv}\icorrload{fj}{fig1j_meta.csv}%
\pgfplotstableread[col sep=comma]{\icorrtikzdata fig1_curves.csv}\mdcurves
\pgfplotstableread[col sep=comma]{\icorrtikzdata fig1j_pair.csv}\mdpair
\pgfplotstableread[col sep=comma]{\icorrtikzdata fig1_prior.csv}\mdprior
\pgfplotstableread[col sep=comma]{\icorrtikzdata fig1_fan_k0.csv}\mdfanA
\pgfplotstableread[col sep=comma]{\icorrtikzdata fig1_fan_k6.csv}\mdfanB
\pgfplotstableread[col sep=comma]{\icorrtikzdata fig1_fan_k24.csv}\mdfanC
\pgfplotstableread[col sep=comma]{\icorrtikzdata fig1_sigmoid.csv}\mdsigmoid
\pgfplotstablegetcolsof{\mdfanA}\pgfmathtruncatemacro\mdlastdraw{\pgfplotsretval-2}% d0..dN
\pgfplotstablegetcolsof{\mdprior}\pgfmathtruncatemacro\mdlastprior{\pgfplotsretval-3}% s0..sN
\tikzset{
  % a line of text centred on the given y (its baseline 0.257 em below)
  md text/.style={inner sep=0pt, icInk, anchor=base west, yshift=-0.257*7pt*\mds},
  md head/.style={inner sep=0pt, icInk, anchor=base west, yshift=-0.257*8pt*\mds,
                  font=\footnotesize\bfseries},
  % several lines, centred on y
  md block/.style={inner sep=0pt, icInk, anchor=west, align=left,
                   font=\fontsize{7}{7.7}\selectfont},
  md box/.style={draw=icMuted, fill=white, line width=\mds*0.6pt, rounded corners=0.05*\mdu},
  md arrow/.style={draw=#1, line width=\mds*0.8pt,
                   -{Triangle[length=\mds*2.8pt, width=\mds*2.8pt]}},
  md arrow/.default=icMuted,
  md generic/.style={icGeneric, line width=\mds*1.0pt},
}
\pgfplotsset{
  md mini/.style={anchor=south west, scale only axis, hide axis, clip=false,
    xmin=\icv{fm}{0}{xmin}, xmax=\icv{fm}{0}{xmax},
    ymin=\icv{fm}{0}{ymin}, ymax=\icv{fm}{0}{ymax}},
  md line/.style={line cap=round, line join=round, mark=none},
}
% a small bare axes: gait cycle on x, torque on y, a zero line; #1, #2 lower-left corner,
% #3, #4 width and height (in), #5 its curves; the style md extra: further axis options
\pgfplotsset{md extra/.style={}}
\def\mdmini#1#2#3#4#5{%
  \begin{axis}[md mini, at={({(#1)*\mdu},{(#2)*\mdu})}, width=(#3)*\mdu, height=(#4)*\mdu,
               md extra]
    \draw[icRule, line width=\mds*0.5pt]
      (axis cs:\icv{fm}{0}{xmin},0) -- (axis cs:\icv{fm}{0}{xmax},0);
    #5
    \draw[icContext, line width=\mds*0.6pt]
      (rel axis cs:0,1) -- (rel axis cs:0,0) -- (rel axis cs:1,0);
  \end{axis}}
% a fan of posterior particles (#1 the table, #2 colour, #3 opacity)
\def\mdfan#1#2#3{%
  \pgfplotsinvokeforeach{0,...,\mdlastdraw}{%
    \addplot[md line, #2, line width=\mds*0.4pt, draw opacity=#3] table[x=phase, y=d##1] {#1};}}
% the prior's shape draws
\def\mdpriorfan{%
  \pgfplotsinvokeforeach{0,...,\mdlastprior}{%
    \addplot[md line, icContext, line width=\mds*0.35pt, draw opacity=0.7]
      table[x=phase, y=s##1] {\mdprior};}}
\def\mdpreferred#1{\addplot[md line, icInk, line width=\mds*#1pt] table[x=phase, y=preferred]
  {\mdcurves};}
\def\mdestimate#1#2{\addplot[md line, icBlue, line width=\mds*#2pt] table[x=phase, y=#1]
  {\mdcurves};}
\def\mdgeneric{\addplot[md generic, md line] table[x=phase, y=generic] {\mdcurves};}
% a profile of the pair (#1 its column, #2 colour), the preferred profile, the gap between them
\def\mdgap#1#2{%
  \addplot[name path=mdgapA, draw=none, mark=none] table[x=phase, y=#1] {\mdpair};
  \addplot[name path=mdgapB, draw=none, mark=none] table[x=phase, y=preferred] {\mdpair};
  \addplot[#2, fill opacity=0.25, draw=none] fill between[of=mdgapA and mdgapB];
  \addplot[md line, icInk, line width=\mds*0.8pt] table[x=phase, y=preferred] {\mdpair};
  \addplot[md line, #2, line width=\mds*1.1pt] table[x=phase, y=#1] {\mdpair};}
% columns and rows (in)
\def\LXa{0.03}\def\LXb{0.77}\def\MXa{0.86}\def\MXb{2.58}\def\RXa{2.67}\def\RXb{3.47}
\def\TYa{0.80}\def\TYb{2.46}
\pgfmathsetmacro\lc{(\LXa+\LXb)/2}\pgfmathsetmacro\rc{(\RXa+\RXb)/2}
% ---- legend: one centred row (0.18 in lines, 0.05 in pad, 0.16 in gaps), widths of the labels
% measured in this type
\pgfmathsetmacro\lwa{width("PsyPref")/68.4}
\pgfmathsetmacro\lwb{width("preferred profile")/68.4}
\pgfmathsetmacro\lwc{width("generic assistance")/68.4}
\pgfmathsetmacro\lxa{(3.5-(3*(0.18+0.05)+\lwa+\lwb+\lwc+2*0.16))/2}
\pgfmathsetmacro\lxb{\lxa+0.18+0.05+\lwa+0.16}
\pgfmathsetmacro\lxc{\lxb+0.18+0.05+\lwb+0.16}
\def\LGy{2.548}
\draw[icBlue, line width=\mds*1.2pt] (\lxa,\LGy) -- ++(0.18,0);
\node[md text] at (\lxa+0.23,\LGy) {PsyPref};
\draw[icInk, line width=\mds*0.9pt] (\lxb,\LGy) -- ++(0.18,0);
\node[md text] at (\lxb+0.23,\LGy) {Preferred profile};
\draw[md generic] (\lxc,\LGy) -- ++(0.18,0);
\node[md text] at (\lxc+0.23,\LGy) {Generic assistance};
% ---- left: what the prior is built from <- the study's peak-torque range; italics: where prior
% knowledge enters
\draw[md box] (\LXa,2.07) rectangle (\LXb,\TYb);
\node[md block, anchor=center, align=center] at (\lc,2.19) {Walking data\\\textit{or other curves}};
\draw[md arrow] (\lc,2.065) -- (\lc,1.985);
\draw[md box] (\LXa,1.25) rectangle (\LXb,1.98);
\node[md block, anchor=center, align=center] at (\lc,1.835) {Profile-shape\\prior, \textit{width}};
\node[md text, anchor=base] at (\LXa+0.15+0.27,1.315) {Gait cycle};
\node[inner sep=0pt, icInk, rotate=90, anchor=center] at (\LXa+0.075,1.40+0.14) {Torque};
\draw[icMuted, line width=\mds*0.6pt, {Triangle[length=\mds*2.2pt, width=\mds*2.2pt]}-{Triangle[length=\mds*2.2pt, width=\mds*2.2pt]}]
  (\LXb-0.033,1.42) -- (\LXb-0.033,1.68);
\draw[md arrow] (\lc,1.195) -- (\lc,1.245);
\draw[md box] (\LXa,\TYa) rectangle (\LXb,1.19);
\node[md block, anchor=center, align=center] at (\lc,0.995)
  {Peak torque\\range or level\\\textit{device, protocol}};
\draw[md arrow] (\LXb+0.005,1.89) -- (\MXa-0.005,1.89);
% ---- middle top: PsyPref's belief
\draw[md box, fill=icTintBlue, draw=icBlue, line width=\mds*0.9pt] (\MXa,1.70) rectangle (\MXb,\TYb);
\node[md head, text=icBlue] at (\MXa+0.07,2.375) {PsyPref:};
\node[md text] at (\MXa+0.07,2.258) {belief over the preferred profile $\mathbf z^\ast$,};
\node[md text] at (\MXa+0.07,2.146) {discrimination scale $\ell$, lapse rate $\lambda$};
\pgfmathsetmacro\pxa{(\MXa+\MXb)/2-(3*0.46+2*0.115)/2}
\pgfmathsetmacro\pxb{\pxa+0.46+0.115}\pgfmathsetmacro\pxc{\pxb+0.46+0.115}
\def\SNy{1.835}\def\SNh{0.24}% the belief insets: bottom and height
\node[md text, anchor=base] at (\pxa+0.23,1.77) {$k = \icv{fm}{0}{k0}$};
\node[md text, anchor=base] at (\pxb+0.23,1.77) {$k = \icv{fm}{0}{k1}$};
\node[md text, anchor=base] at (\pxc+0.23,1.77) {$k = \icv{fm}{0}{k2}$};
\draw[md arrow=icContext, line width=\mds*0.7pt]
  (\pxb-0.115+0.025,\SNy+\SNh/2) -- (\pxb-0.025,\SNy+\SNh/2);
\draw[md arrow=icContext, line width=\mds*0.7pt]
  (\pxc-0.115+0.025,\SNy+\SNh/2) -- (\pxc-0.025,\SNy+\SNh/2);
% ---- middle bottom: the user
\draw[md box, fill=icTintBlue] (\MXa,\TYa) rectangle (\MXb,1.40);
\pgfmathsetmacro\qx{\MXa+0.06+0.44+0.08}
\node[md block, font=\fontsize{7}{8}\selectfont\bfseries] at (\qx,1.27)
  {Profile A, B or\\no preference?};
% ---- the loop's two edges
\pgfmathsetmacro\xl{\MXa+0.10}\pgfmathsetmacro\xr{\MXb-0.10}
\draw[md arrow=icInk] (\xr,1.695) -- (\xr,1.405);
\node[md block, anchor=east, align=right] at (\xr-0.05,1.55) {Most informative\\pair};
\draw[md arrow=icInk] (\xl,1.405) -- (\xl,1.695);
\node[md block] at (\xl+0.05,1.55) {Belief\\update};
% ---- right: identified assistance -> validation
\draw[md arrow] (\MXb+0.005,2.01) -- (\RXa-0.005,2.01);
\draw[md box] (\RXa,1.70) rectangle (\RXb,\TYb);
\node[md block] at (\RXa+0.06,2.2024)
  {\textcolor{icBlue}{Identified}\\\textcolor{icBlue}{assistance}\\after $k$\\comparisons};
\draw[md arrow] (\rc,1.695) -- (\rc,1.615);
\draw[md box] (\RXa,\TYa) rectangle (\RXb,1.61);
\node[md text, font=\scriptsize\bfseries] at (\RXa+0.06,1.515) {Validation:};
\node[md block] at (\RXa+0.06,1.30) {identified vs.\\generic\\assistance};
% ---- bottom: this paper's simulated users
% the worked answer: the plots' bottom and height, the pair's left edges and width, the answer
% model's left edge and width, the row of labels below them, how far the box grows downward (in)
\def\WEy{0.07}\def\WEh{0.34}\def\WEa{1.07}\def\WEb{1.68}\def\WEw{0.52}
\def\SGx{2.38}\def\SGw{0.54}\def\SGy{\WEy}\def\SGh{\WEh}
\def\WEl{-0.025}\def\WEd{0.12}
\draw[draw=icMuted, line width=\mds*0.8pt, dash pattern=on \mds*2.4pt off \mds*1.6pt,
      rounded corners=0.06*\mdu] (\LXa,0.03-\WEd) rectangle (\RXb,0.70);
\node[md head] at (\LXa+0.07,0.615) {This paper: simulated users};
% who the users are; how they answer: the pair's distances, then the probability they give
\node[md block, anchor=north west] at (\LXa+0.07,0.535)
  {\textbf{Preferred profile}\\\textit{Shape:} walker not\\\phantom{\textit{Shape:} }in the prior\\\textit{Peaks:} study's range};
\node[md block, anchor=north west] at (\WEa,0.535) {\textbf{Distance} $d$ (Nm)};
\node[md block, anchor=north west] at (2.12,0.535) {\textbf{Answer model:} $P$(choose A)};
\node[md text, anchor=base] at (\WEa+\WEw/2,\WEl) {$d(\mathrm{A})\!=\!\icv{fj}{0}{d1_txt}$};
\node[md text, anchor=base] at (\WEb+\WEw/2,\WEl) {$d(\mathrm{B})\!=\!\icv{fj}{0}{d2_txt}$};
\node[md text, anchor=base] at (\SGx+\SGw/2,\WEl) {$d(\mathrm{B}) - d(\mathrm{A})$};
\draw[md arrow] (\WEb+\WEw+0.03,\WEy+\WEh/2) -- (\SGx-0.05,\WEy+\WEh/2);
\draw[md arrow] (\MXa+0.28,0.705) -- (\MXa+0.28,0.795);
% ---- the walkers: three without, the user with the hip exoskeleton
\icfwalker{(\LXa+0.20)*\mdu}{2.315*\mdu}{0.12*\mdu}{5}
\icfwalker{(\LXa+0.38)*\mdu}{2.315*\mdu}{0.12*\mdu}{5}
\icfwalker{(\LXa+0.56)*\mdu}{2.315*\mdu}{0.12*\mdu}{5}
\icfwalker{(\MXa+0.06+0.22-0.02)*\mdu}{(\TYa+0.05+0.02)*\mdu}{0.46*\mdu}{11}
% ---- the small axes
% the prior's shapes at unit peaks: limits of their own
\pgfplotsset{md extra/.style={ymin=\icv{fm}{0}{prior_ymin}, ymax=\icv{fm}{0}{prior_ymax}}}
\mdmini{\LXa+0.15}{1.40}{0.54}{0.29}{%
  \mdpriorfan\addplot[md generic, md line] table[x=phase, y=generic_shape] {\mdprior};}
\pgfplotsset{md extra/.style={}}
\mdmini{\pxa}{\SNy}{0.46}{\SNh}{%
  \mdfan{\mdfanA}{icBlueLight}{0.55}\mdestimate{mean_k0}{1.0}\mdpreferred{0.8}}
\mdmini{\pxb}{\SNy}{0.46}{\SNh}{%
  \mdfan{\mdfanB}{icBlueLight}{0.55}\mdestimate{mean_k6}{1.0}\mdpreferred{0.8}}
\mdmini{\pxc}{\SNy}{0.46}{\SNh}{%
  \mdfan{\mdfanC}{icBlueLight}{0.55}\mdestimate{mean_k24}{1.0}\mdpreferred{0.8}}
% the pair: each label inside its plot's top-left corner, clear of the other plot
\mdmini{\qx}{0.87}{0.46}{0.27}{%
  \addplot[md line, icPairA, line width=\mds*1.2pt] table[x=phase, y=pair1] {\mdpair};
  \node[inner sep=0pt, anchor=north west, font=\scriptsize\bfseries, icInk, xshift=\mds*2.5pt,
        yshift=\mds*1.5pt] at (rel axis cs:0,1) {A};}
\mdmini{\qx+0.60}{0.87}{0.46}{0.27}{%
  \addplot[md line, icPairB, line width=\mds*1.2pt] table[x=phase, y=pair2] {\mdpair};
  \node[inner sep=0pt, anchor=north west, font=\scriptsize\bfseries, icInk, xshift=\mds*2.5pt,
        yshift=\mds*1.5pt] at (rel axis cs:0,1) {B};}
% the same pair with the preferred profile: limits of its own, so that the gaps fill the plots
\pgfplotsset{md extra/.style={ymin=\icv{fj}{0}{ymin}, ymax=\icv{fj}{0}{ymax}}}
\mdmini{\WEa}{\WEy}{\WEw}{\WEh}{%
  \mdgap{pair1}{icPairA}
  \node[inner sep=0pt, anchor=north west, font=\scriptsize\bfseries, icInk, xshift=\mds*2.5pt,
        yshift=\mds*1.5pt] at (rel axis cs:0,1) {A};}
\mdmini{\WEb}{\WEy}{\WEw}{\WEh}{%
  \mdgap{pair2}{icPairB}
  \node[inner sep=0pt, anchor=north west, font=\scriptsize\bfseries, icInk, xshift=\mds*2.5pt,
        yshift=\mds*1.5pt] at (rel axis cs:0,1) {B};}
\pgfplotsset{md extra/.style={}}
\mdmini{\RXa+0.06}{1.75}{\RXb-\RXa-0.12}{0.21}{%
  \mdpreferred{0.9}\mdestimate{mean_k24}{1.2}
  \node[inner sep=0pt, anchor=south east, icMuted]
    at (rel axis cs:0.99,0.02) {$k = \icv{fm}{0}{k_id}$};}
\mdmini{\RXa+0.06}{\TYa+0.05}{\RXb-\RXa-0.12}{0.27}{%
  \mdgeneric\mdestimate{mean_k24}{1.2}}
% ---- the answer model: lapse floor and ceiling, ticks on the right
\pgfmathsetmacro\mdlo{\icv{fm}{0}{lapse}/2}\pgfmathsetmacro\mdhi{1-\icv{fm}{0}{lapse}/2}
\begin{axis}[at={(\SGx*\mdu,\SGy*\mdu)}, anchor=south west, scale only axis,
  width=\SGw*\mdu, height=\SGh*\mdu, clip=false,
  xmin=\icv{fm}{0}{sig_xmin}, xmax=\icv{fm}{0}{sig_xmax},
  ymin=\icv{fm}{0}{sig_ymin}, ymax=\icv{fm}{0}{sig_ymax},
  axis x line*=bottom, axis y line*=right, axis line style={-, icMuted, line width=\mds*0.6pt},
  tick align=outside, major tick length=\mds*2pt, tick style={icMuted, line width=\mds*0.6pt},
  xtick={0}, xticklabels={}, ytick={\mdlo,0.5,\mdhi},
  yticklabels={$\lambda$/2, 0.5, 1\,$-$\,$\lambda$/2}, yticklabel pos=right,
  yticklabel style={icInk, font=\scriptsize, inner sep=0pt, xshift=\mds*1.5pt},
]
\draw[icContext, line width=\mds*0.6pt, dash pattern=on \mds*0.6pt off \mds*0.9pt]
  (axis cs:\icv{fm}{0}{sig_xmin},\mdlo) -- (axis cs:\icv{fm}{0}{sig_xmax},\mdlo);
\draw[icContext, line width=\mds*0.6pt, dash pattern=on \mds*0.6pt off \mds*0.9pt]
  (axis cs:\icv{fm}{0}{sig_xmin},\mdhi) -- (axis cs:\icv{fm}{0}{sig_xmax},\mdhi);
\draw[icRule, line width=\mds*0.5pt]
  (axis cs:0,\icv{fm}{0}{sig_ymin}) -- (axis cs:0,\icv{fm}{0}{sig_ymax});
\addplot[icInk, line width=\mds*1.2pt, mark=none, line join=round] table[x=delta, y=p]
  {\mdsigmoid};
% the pair's reading: its distance difference, the probability it gives
\draw[icMuted, line width=\mds*0.6pt, dash pattern=on \mds*0.9pt off \mds*0.9pt]
  (axis cs:\icv{fj}{0}{delta},\icv{fm}{0}{sig_ymin}) -- (axis cs:\icv{fj}{0}{delta},\icv{fj}{0}{p})
  -- (axis cs:\icv{fm}{0}{sig_xmin},\icv{fj}{0}{p});
\fill[icInk] (axis cs:\icv{fj}{0}{delta},\icv{fj}{0}{p}) circle[radius=\mds*1.3pt];
\node[inner sep=0pt, anchor=north west, font=\scriptsize\bfseries, icInk, yshift=-\mds*1.2pt]
  at (axis cs:\icv{fm}{0}{sig_xmin},\icv{fj}{0}{p}) {\icv{fj}{0}{p_txt}};
\end{axis}
\pgfresetboundingbox
\useasboundingbox (0,-\WEd) rectangle (3.5,2.6);
\end{tikzpicture}%
}
\end{document}
"
% The dark colors follow the README's other figures (THEMES["dark"] in scripts/make_readme_assets.py).
\documentclass[conference]{IEEEtran}
\usepackage{amsmath,amssymb}
\usepackage{xcolor}
\usepackage{pgfplots}
\pgfplotsset{compat=1.18}
\usepgfplotslibrary{fillbetween}
\usetikzlibrary{arrows.meta}
% icorr_style.tex: shared colours and line styles of the amplitude-free paper's figures
% (paper/notes/story_concept.md section 7a, scheme C). \input once in the preamble, after
% \usepackage{pgfplots}. Every figure uses these names; no figure defines its own colours.
%
% Roles:
%   PsyPref (all variants) ............ icBlue, solid
%   best published comparator ......... icOrange, solid, hollow marker: square CMA-ES, triangle Arens' GP
%   GP-BALD (same prior) .............. icGrey, dashed
%   CMA-ES on Schaefer's family ....... icDark, dash-dot
%   generic assistance (50 %) ......... icGeneric, solid thin
%   best K-weight profile (ceiling, floor) icContext, dotted
%   K = 8 / 16 / 24 (or k = 24 / 40 / 150) icBlueLight / icBlue / icBlueDark
%   settings as markers ............... circle 3-35 Nm, square fixed 12 Nm, diamond 3-12 Nm
%   robustness: filled = fixed yardstick (post hoc), hollow = registered scoring
\definecolor{icBlue}{HTML}{1845FB}
\definecolor{icBlueLight}{HTML}{8FA4FD}
\definecolor{icBlueDark}{HTML}{0F2C9E}
\definecolor{icOrange}{HTML}{F89C20}
\definecolor{icOrangeText}{HTML}{B86E00}% text in the comparator's hue (contrast 4.0 on white)
\definecolor{icGrey}{HTML}{7A7A7A}
\definecolor{icDark}{HTML}{333333}
\definecolor{icGeneric}{HTML}{9C9CA1}
\definecolor{icContext}{HTML}{BBBBBB}
\definecolor{icRule}{HTML}{DDDDDD}
\definecolor{icInk}{HTML}{222222}
\definecolor{icMuted}{HTML}{666666}
% Fig. 1 only: the pair shown to the user (author's choice, 2026-10-29: distinct hues read
% stronger than two blues; gold was dropped because it equals the comparator's orange)
\definecolor{icTintBlue}{HTML}{F1F3F8}% Fig. 1: the loop, PsyPref's and the user's box (author's choice)
\definecolor{icPairA}{HTML}{BB5566}
\definecolor{icPairB}{HTML}{33BBEE}
\tikzset{
  ic psypref/.style={icBlue, line width=1.5pt, line cap=round, line join=round},
  ic best/.style={icOrange, line width=1.2pt, line cap=round, line join=round},
  ic gpbald/.style={icGrey, line width=0.9pt, dash pattern=on 3.2pt off 1.6pt},
  ic sch/.style={icDark, line width=0.9pt, dash pattern=on 3.6pt off 1.4pt on 0.9pt off 1.4pt},
  ic generic/.style={icGeneric, line width=0.6pt},
  ic context/.style={icContext, line width=0.8pt, dash pattern=on 0.9pt off 1.5pt},
  ic rule/.style={icRule, line width=0.5pt},
  ic label/.style={inner sep=0.6pt, font=\scriptsize, align=left},
  ic note/.style={inner sep=0.8pt, font=\scriptsize, icInk, fill=white, fill opacity=0.85,
                  text opacity=1, align=right},
  ic band/.style={fill=#1, fill opacity=0.15, draw=none},
}
\pgfplotsset{
  ic axis/.style={
    axis line style={icMuted, line width=0.5pt},
    tick style={icMuted, line width=0.5pt},
    tick label style={font=\scriptsize, icMuted},
    label style={font=\footnotesize, icInk},
    title style={font=\footnotesize, icInk},
    axis x line*=bottom, axis y line*=left,
    tick align=outside, major tick length=2pt,
    clip=false,
  },
}

\ifdefined\readmedark
  \definecolor{icInk}{HTML}{E6EDF3}
  \definecolor{icMuted}{HTML}{9DA7B3}
  \definecolor{icRule}{HTML}{30363D}
  \definecolor{icBlue}{HTML}{8FA4FD}
  \definecolor{icBlueLight}{HTML}{3F57C8}
  \definecolor{icBlueDark}{HTML}{C3CEFF}
  \definecolor{icGeneric}{HTML}{8B949E}
  \definecolor{icContext}{HTML}{4B535D}
  \definecolor{icGrey}{HTML}{9DA7B3}
  \definecolor{icDark}{HTML}{C9D1D9}
  \definecolor{icPairA}{HTML}{E58A9A}
  \definecolor{icPairB}{HTML}{33BBEE}
  \definecolor{icOrangeText}{HTML}{F8B860}
  \definecolor{icTintBlue}{HTML}{161D33}% PsyPref's and the user's box
  \definecolor{white}{HTML}{161B22}% the other boxes, slightly raised from GitHub's dark background
\fi
\newcommand{\icorrtikzdata}{data/}
% fig3p_text.tex: written by scripts/paper_v2_fig3p.py (repository); do not edit by hand.
% Every number of Fig. 3's caption and the draft Results paragraph, formatted from values the script checked
% (fam_report.md, fam_numbers.json, amp_table1.json, Fig. 2's data). Differences as fam_report.md prints them.
\def\pjK{24}
\def\pjNusers{200}
\def\pjNsubj{5}
\def\pjSigma{1}
\def\pjLapse{0.1}
\def\pjGapW{18.8}
\def\pjGapWci{(14.7, 22.8)}
\def\pjGapF{9.0}
\def\pjGapFci{(6.1, 11.9)}
\def\pjSchCeilW{94.1}
\def\pjSchCeilF{88.9}
\def\pjSchPsyW{86.4}
\def\pjSchPsyF{74.8}
\def\pjCmaW{67.6}
\def\pjCmaWci{(63.1, 71.9)}
\def\pjCmaF{65.8}
\def\pjCmaFci{(63.5, 67.9)}
\def\pjPubSch{76.6}
\def\pjPubSchci{(74.6, 78.5)}
\def\pjPubPchip{86.6}
\def\pjPubPchipci{(85.0, 88.3)}
\def\pjPubSchDiff{17.5}
\def\pjPubSchDiffci{(15.6, 19.5)}
\def\pjPubPchipDiff{7.4}
\def\pjPubPchipDiffci{(5.9, 9.0)}
\def\pjProjSchW{6.1}
\def\pjProjSchWci{($-7.8$, $-4.6$)}
\def\pjProjSchWsigned{$-6.1$}
\def\pjProjFranksW{$-0.6$}
\def\pjProjPchipW{$-1.1$}
\def\pjProjFixLo{3.4}
\def\pjProjFixHi{9.8}
\def\pjPrompPsyW{92.5}
\def\pjPubTimeFranksW{91.8}
\def\pjPubTimeFranksF{74.2}
\def\pjPubTimePchipW{93.9}
\def\pjPubTimePchipF{89.1}
\def\pjExKpair{15}

% fig4m_text.tex: written by scripts/paper_v2_fig4m.py (repository); do not edit by hand.
\def\fmSThreeW{$+16.4$ ($+14.4$, $+18.3$)}
\def\fmSThreeF{$+7.9$ ($+6.1$, $+9.2$)}
\def\fmSThreeLevW{75.3}
\def\fmSThreeLevF{63.6}
\def\fmStartC{1.5}
\def\fmStartP{2.8}
\def\fmK{24}
\def\fmNclear{68}
\def\fmNall{75}
\def\fmNcell{70}
\def\fmNfix{68}
\def\fmNusers{200}
\def\fmNother{120}
\def\fmNtie{160}
\def\fmMainW{$+13.1$ ($+10.3$, $+14.7$)}
\def\fmMainF{$+8.8$ ($+6.8$, $+10.5$)}
\def\fmSingleW{$+2.6$ ($-0.5$, $+4.8$)}
\def\fmSingleF{$+3.4$ ($+0.5$, $+6.2$)}
\def\fmOtherW{$+16.0$ ($+12.6$, $+19.2$)}
\def\fmOtherF{$+12.0$ ($+8.5$, $+14.3$)}
\def\fmPeaks{$+8.1$}
\def\fmLaterW{$+12.1$}
\def\fmLaterF{$+10.7$}
\def\fmNoneF{$-2.0$ ($-5.3$, $+1.0$)}
\def\fmShapesF{$-3.1$ ($-4.8$, $-1.4$)}
\def\fmShapesW{$+13.7$ ($+10.8$, $+15.7$)}
\def\fmShapesNext{$+15.6$}
\def\fmOtherShapesW{$+5.0$}
\def\fmOtherShapesF{$+5.1$}
\def\fmPessW{$+12.6$ ($+10.0$, $+14.6$)}
\def\fmPessF{$-1.3$ ($-2.9$, $+0.3$)}
\def\fmDropSingle{17.2}
\def\fmDropPess{21.9}
\def\fmFixGap{8.6}
\def\fmBaldLo{1.9}
\def\fmBaldHi{5.7}

\pagestyle{empty}
\begin{document}
\thispagestyle{empty}
\noindent{\color{icInk}% fig1.tex: Fig. 1 of the amplitude-free paper (the method diagram), in TikZ: 0.95 x 252 = 239.4 pt
% wide and 177.84 pt + \WEd in = 186.05 pt tall, for
%   \input{figures/tikz_v2/fig1.tex}
% The simulated users' box shows a worked answer: the pair the loop asks, its distances to the
% preferred profile and the probability they give.
% Preamble: the paper's (pgfplots, compat 1.18, arrows.meta, fillbetween) plus
% % icorr_style.tex: shared colours and line styles of the amplitude-free paper's figures
% (paper/notes/story_concept.md section 7a, scheme C). \input once in the preamble, after
% \usepackage{pgfplots}. Every figure uses these names; no figure defines its own colours.
%
% Roles:
%   PsyPref (all variants) ............ icBlue, solid
%   best published comparator ......... icOrange, solid, hollow marker: square CMA-ES, triangle Arens' GP
%   GP-BALD (same prior) .............. icGrey, dashed
%   CMA-ES on Schaefer's family ....... icDark, dash-dot
%   generic assistance (50 %) ......... icGeneric, solid thin
%   best K-weight profile (ceiling, floor) icContext, dotted
%   K = 8 / 16 / 24 (or k = 24 / 40 / 150) icBlueLight / icBlue / icBlueDark
%   settings as markers ............... circle 3-35 Nm, square fixed 12 Nm, diamond 3-12 Nm
%   robustness: filled = fixed yardstick (post hoc), hollow = registered scoring
\definecolor{icBlue}{HTML}{1845FB}
\definecolor{icBlueLight}{HTML}{8FA4FD}
\definecolor{icBlueDark}{HTML}{0F2C9E}
\definecolor{icOrange}{HTML}{F89C20}
\definecolor{icOrangeText}{HTML}{B86E00}% text in the comparator's hue (contrast 4.0 on white)
\definecolor{icGrey}{HTML}{7A7A7A}
\definecolor{icDark}{HTML}{333333}
\definecolor{icGeneric}{HTML}{9C9CA1}
\definecolor{icContext}{HTML}{BBBBBB}
\definecolor{icRule}{HTML}{DDDDDD}
\definecolor{icInk}{HTML}{222222}
\definecolor{icMuted}{HTML}{666666}
% Fig. 1 only: the pair shown to the user (author's choice, 2026-10-29: distinct hues read
% stronger than two blues; gold was dropped because it equals the comparator's orange)
\definecolor{icTintBlue}{HTML}{F1F3F8}% Fig. 1: the loop, PsyPref's and the user's box (author's choice)
\definecolor{icPairA}{HTML}{BB5566}
\definecolor{icPairB}{HTML}{33BBEE}
\tikzset{
  ic psypref/.style={icBlue, line width=1.5pt, line cap=round, line join=round},
  ic best/.style={icOrange, line width=1.2pt, line cap=round, line join=round},
  ic gpbald/.style={icGrey, line width=0.9pt, dash pattern=on 3.2pt off 1.6pt},
  ic sch/.style={icDark, line width=0.9pt, dash pattern=on 3.6pt off 1.4pt on 0.9pt off 1.4pt},
  ic generic/.style={icGeneric, line width=0.6pt},
  ic context/.style={icContext, line width=0.8pt, dash pattern=on 0.9pt off 1.5pt},
  ic rule/.style={icRule, line width=0.5pt},
  ic label/.style={inner sep=0.6pt, font=\scriptsize, align=left},
  ic note/.style={inner sep=0.8pt, font=\scriptsize, icInk, fill=white, fill opacity=0.85,
                  text opacity=1, align=right},
  ic band/.style={fill=#1, fill opacity=0.15, draw=none},
}
\pgfplotsset{
  ic axis/.style={
    axis line style={icMuted, line width=0.5pt},
    tick style={icMuted, line width=0.5pt},
    tick label style={font=\scriptsize, icMuted},
    label style={font=\footnotesize, icInk},
    title style={font=\footnotesize, icInk},
    axis x line*=bottom, axis y line*=left,
    tick align=outside, major tick length=2pt,
    clip=false,
  },
}
; colours are icorr_style.tex's names only.
% Data: \icorrtikzdata (default figures/tikz_v2/data/), fig1_*.csv (scripts/paper_v2_fig1.py,
% repository): one PsyPref-T session in the wide setting (3-35 Nm) for the lower of the two
% middle users of FULL's 200 by distance after 24 comparisons (rank 100 of 200), rerun as FULL's
% run_one builds it and checked against FULL's stored distances; belief particles and the delivered
% profile at k = 0, 6, 24, the preferred profile, generic assistance, 30 draws of the walking-data
% shape prior (unit peaks), the answer model's sigmoid and the walker silhouette. fig1j_*.csv
% (scripts/paper_v2_fig1j.py, the same session): the pair, by rule the one of the first 24
% comparisons with d(A) < d(B) whose d(B) - d(A) lies closest to 1 JND; its distances and
% P(choose A). The CSV columns keep pair1 and pair2 for A and B, and fig1_sigmoid.csv's p is the
% probability of the first profile.
% Layout: the old Fig. 1's coordinates in inches of the script (3.5 x 2.6 in), scaled as the
% paper scales it: 1 in = 0.95 x 72 pt = \mdu; line widths scaled alike (\mds).
\providecommand{\icorrtikzdata}{figures/tikz_v2/data/}%
\makeatletter
\providecommand{\icorrload}[2]{% \icorrload{name}{file.csv}: every cell as \icv{name}{row}{column}%
  \pgfplotstableread[col sep=comma]{\icorrtikzdata #2}\ic@table
  \pgfplotstableforeachcolumn\ic@table\as\ic@col{%
    \pgfplotstableforeachcolumnelement{\ic@col}\of\ic@table\as\ic@cell{%
      \expandafter\global\expandafter\let\csname ic@#1@\pgfplotstablerow @\ic@col\endcsname
        \ic@cell}}}%
\providecommand{\icv}[3]{\csname ic@#1@#2@#3\endcsname}%
% the walker: \icfwalker{x}{y}{height}{last part}: hip at x, ground at y, height, all three
% lengths; last part 5 draws the walker alone, 11 with the hip exoskeleton. Geometry from
% fig1_walker_part<i>.csv (units of the height), style and colour names from
% fig1_walker_parts.csv.
\providecommand{\icfwalker}[4]{%
  \begin{axis}[at={({#1},{#2})}, anchor=origin, x=#3, y=#3, hide axis, clip=false,%
               xmin=-0.3, xmax=0.3, ymin=0, ymax=1]%
    \pgfplotsinvokeforeach{0,...,#4}{\icf@walkerpart{##1}{#3}}%
  \end{axis}}%
\def\icf@none{none}\def\icf@line{line}%
\providecommand{\icf@walkerpart}[2]{%
  \edef\icf@edge{\icv{fwk}{#1}{edge}}\edef\icf@fill{\icv{fwk}{#1}{fill}}%
  \edef\icf@kind{\icv{fwk}{#1}{kind}}%
  \ifx\icf@edge\icf@none\else\colorlet{icfwe#1}{\icf@edge}\fi
  \ifx\icf@fill\icf@none\else\colorlet{icfwf#1}{\icf@fill}\fi
  \ifx\icf@kind\icf@line
    \addplot[draw=icfwe#1, line width=\icv{fwk}{#1}{lw_per_height}*#2,%
             line cap=\icv{fwk}{#1}{cap}, line join=round, mark=none]%
      table[col sep=comma, x=x, y=y] {\icorrtikzdata fig1_walker_part#1.csv};%
  \else\ifx\icf@edge\icf@none
    \addplot[fill=icfwf#1, draw=none, mark=none]%
      table[col sep=comma, x=x, y=y] {\icorrtikzdata fig1_walker_part#1.csv};%
  \else
    \addplot[fill=icfwf#1, draw=icfwe#1, line width=\icv{fwk}{#1}{lw_per_height}*#2, mark=none]%
      table[col sep=comma, x=x, y=y] {\icorrtikzdata fig1_walker_part#1.csv};%
  \fi\fi}%
\makeatother
\begin{tikzpicture}[x=68.4pt, y=68.4pt, font=\scriptsize]
\scriptsize% the type the legend row is measured in
\def\mdu{68.4pt}% 1 in of the script, as printed
\def\mds{0.95}% the paper's scale: line widths
\icorrload{fm}{fig1_meta.csv}\icorrload{fwk}{fig1_walker_parts.csv}\icorrload{fj}{fig1j_meta.csv}%
\pgfplotstableread[col sep=comma]{\icorrtikzdata fig1_curves.csv}\mdcurves
\pgfplotstableread[col sep=comma]{\icorrtikzdata fig1j_pair.csv}\mdpair
\pgfplotstableread[col sep=comma]{\icorrtikzdata fig1_prior.csv}\mdprior
\pgfplotstableread[col sep=comma]{\icorrtikzdata fig1_fan_k0.csv}\mdfanA
\pgfplotstableread[col sep=comma]{\icorrtikzdata fig1_fan_k6.csv}\mdfanB
\pgfplotstableread[col sep=comma]{\icorrtikzdata fig1_fan_k24.csv}\mdfanC
\pgfplotstableread[col sep=comma]{\icorrtikzdata fig1_sigmoid.csv}\mdsigmoid
\pgfplotstablegetcolsof{\mdfanA}\pgfmathtruncatemacro\mdlastdraw{\pgfplotsretval-2}% d0..dN
\pgfplotstablegetcolsof{\mdprior}\pgfmathtruncatemacro\mdlastprior{\pgfplotsretval-3}% s0..sN
\tikzset{
  % a line of text centred on the given y (its baseline 0.257 em below)
  md text/.style={inner sep=0pt, icInk, anchor=base west, yshift=-0.257*7pt*\mds},
  md head/.style={inner sep=0pt, icInk, anchor=base west, yshift=-0.257*8pt*\mds,
                  font=\footnotesize\bfseries},
  % several lines, centred on y
  md block/.style={inner sep=0pt, icInk, anchor=west, align=left,
                   font=\fontsize{7}{7.7}\selectfont},
  md box/.style={draw=icMuted, fill=white, line width=\mds*0.6pt, rounded corners=0.05*\mdu},
  md arrow/.style={draw=#1, line width=\mds*0.8pt,
                   -{Triangle[length=\mds*2.8pt, width=\mds*2.8pt]}},
  md arrow/.default=icMuted,
  md generic/.style={icGeneric, line width=\mds*1.0pt},
}
\pgfplotsset{
  md mini/.style={anchor=south west, scale only axis, hide axis, clip=false,
    xmin=\icv{fm}{0}{xmin}, xmax=\icv{fm}{0}{xmax},
    ymin=\icv{fm}{0}{ymin}, ymax=\icv{fm}{0}{ymax}},
  md line/.style={line cap=round, line join=round, mark=none},
}
% a small bare axes: gait cycle on x, torque on y, a zero line; #1, #2 lower-left corner,
% #3, #4 width and height (in), #5 its curves; the style md extra: further axis options
\pgfplotsset{md extra/.style={}}
\def\mdmini#1#2#3#4#5{%
  \begin{axis}[md mini, at={({(#1)*\mdu},{(#2)*\mdu})}, width=(#3)*\mdu, height=(#4)*\mdu,
               md extra]
    \draw[icRule, line width=\mds*0.5pt]
      (axis cs:\icv{fm}{0}{xmin},0) -- (axis cs:\icv{fm}{0}{xmax},0);
    #5
    \draw[icContext, line width=\mds*0.6pt]
      (rel axis cs:0,1) -- (rel axis cs:0,0) -- (rel axis cs:1,0);
  \end{axis}}
% a fan of posterior particles (#1 the table, #2 colour, #3 opacity)
\def\mdfan#1#2#3{%
  \pgfplotsinvokeforeach{0,...,\mdlastdraw}{%
    \addplot[md line, #2, line width=\mds*0.4pt, draw opacity=#3] table[x=phase, y=d##1] {#1};}}
% the prior's shape draws
\def\mdpriorfan{%
  \pgfplotsinvokeforeach{0,...,\mdlastprior}{%
    \addplot[md line, icContext, line width=\mds*0.35pt, draw opacity=0.7]
      table[x=phase, y=s##1] {\mdprior};}}
\def\mdpreferred#1{\addplot[md line, icInk, line width=\mds*#1pt] table[x=phase, y=preferred]
  {\mdcurves};}
\def\mdestimate#1#2{\addplot[md line, icBlue, line width=\mds*#2pt] table[x=phase, y=#1]
  {\mdcurves};}
\def\mdgeneric{\addplot[md generic, md line] table[x=phase, y=generic] {\mdcurves};}
% a profile of the pair (#1 its column, #2 colour), the preferred profile, the gap between them
\def\mdgap#1#2{%
  \addplot[name path=mdgapA, draw=none, mark=none] table[x=phase, y=#1] {\mdpair};
  \addplot[name path=mdgapB, draw=none, mark=none] table[x=phase, y=preferred] {\mdpair};
  \addplot[#2, fill opacity=0.25, draw=none] fill between[of=mdgapA and mdgapB];
  \addplot[md line, icInk, line width=\mds*0.8pt] table[x=phase, y=preferred] {\mdpair};
  \addplot[md line, #2, line width=\mds*1.1pt] table[x=phase, y=#1] {\mdpair};}
% columns and rows (in)
\def\LXa{0.03}\def\LXb{0.77}\def\MXa{0.86}\def\MXb{2.58}\def\RXa{2.67}\def\RXb{3.47}
\def\TYa{0.80}\def\TYb{2.46}
\pgfmathsetmacro\lc{(\LXa+\LXb)/2}\pgfmathsetmacro\rc{(\RXa+\RXb)/2}
% ---- legend: one centred row (0.18 in lines, 0.05 in pad, 0.16 in gaps), widths of the labels
% measured in this type
\pgfmathsetmacro\lwa{width("PsyPref")/68.4}
\pgfmathsetmacro\lwb{width("preferred profile")/68.4}
\pgfmathsetmacro\lwc{width("generic assistance")/68.4}
\pgfmathsetmacro\lxa{(3.5-(3*(0.18+0.05)+\lwa+\lwb+\lwc+2*0.16))/2}
\pgfmathsetmacro\lxb{\lxa+0.18+0.05+\lwa+0.16}
\pgfmathsetmacro\lxc{\lxb+0.18+0.05+\lwb+0.16}
\def\LGy{2.548}
\draw[icBlue, line width=\mds*1.2pt] (\lxa,\LGy) -- ++(0.18,0);
\node[md text] at (\lxa+0.23,\LGy) {PsyPref};
\draw[icInk, line width=\mds*0.9pt] (\lxb,\LGy) -- ++(0.18,0);
\node[md text] at (\lxb+0.23,\LGy) {Preferred profile};
\draw[md generic] (\lxc,\LGy) -- ++(0.18,0);
\node[md text] at (\lxc+0.23,\LGy) {Generic assistance};
% ---- left: what the prior is built from <- the study's peak-torque range; italics: where prior
% knowledge enters
\draw[md box] (\LXa,2.07) rectangle (\LXb,\TYb);
\node[md block, anchor=center, align=center] at (\lc,2.19) {Walking data\\\textit{or other curves}};
\draw[md arrow] (\lc,2.065) -- (\lc,1.985);
\draw[md box] (\LXa,1.25) rectangle (\LXb,1.98);
\node[md block, anchor=center, align=center] at (\lc,1.835) {Profile-shape\\prior, \textit{width}};
\node[md text, anchor=base] at (\LXa+0.15+0.27,1.315) {Gait cycle};
\node[inner sep=0pt, icInk, rotate=90, anchor=center] at (\LXa+0.075,1.40+0.14) {Torque};
\draw[icMuted, line width=\mds*0.6pt, {Triangle[length=\mds*2.2pt, width=\mds*2.2pt]}-{Triangle[length=\mds*2.2pt, width=\mds*2.2pt]}]
  (\LXb-0.033,1.42) -- (\LXb-0.033,1.68);
\draw[md arrow] (\lc,1.195) -- (\lc,1.245);
\draw[md box] (\LXa,\TYa) rectangle (\LXb,1.19);
\node[md block, anchor=center, align=center] at (\lc,0.995)
  {Peak torque\\range or level\\\textit{device, protocol}};
\draw[md arrow] (\LXb+0.005,1.89) -- (\MXa-0.005,1.89);
% ---- middle top: PsyPref's belief
\draw[md box, fill=icTintBlue, draw=icBlue, line width=\mds*0.9pt] (\MXa,1.70) rectangle (\MXb,\TYb);
\node[md head, text=icBlue] at (\MXa+0.07,2.375) {PsyPref:};
\node[md text] at (\MXa+0.07,2.258) {belief over the preferred profile $\mathbf z^\ast$,};
\node[md text] at (\MXa+0.07,2.146) {discrimination scale $\ell$, lapse rate $\lambda$};
\pgfmathsetmacro\pxa{(\MXa+\MXb)/2-(3*0.46+2*0.115)/2}
\pgfmathsetmacro\pxb{\pxa+0.46+0.115}\pgfmathsetmacro\pxc{\pxb+0.46+0.115}
\def\SNy{1.835}\def\SNh{0.24}% the belief insets: bottom and height
\node[md text, anchor=base] at (\pxa+0.23,1.77) {$k = \icv{fm}{0}{k0}$};
\node[md text, anchor=base] at (\pxb+0.23,1.77) {$k = \icv{fm}{0}{k1}$};
\node[md text, anchor=base] at (\pxc+0.23,1.77) {$k = \icv{fm}{0}{k2}$};
\draw[md arrow=icContext, line width=\mds*0.7pt]
  (\pxb-0.115+0.025,\SNy+\SNh/2) -- (\pxb-0.025,\SNy+\SNh/2);
\draw[md arrow=icContext, line width=\mds*0.7pt]
  (\pxc-0.115+0.025,\SNy+\SNh/2) -- (\pxc-0.025,\SNy+\SNh/2);
% ---- middle bottom: the user
\draw[md box, fill=icTintBlue] (\MXa,\TYa) rectangle (\MXb,1.40);
\pgfmathsetmacro\qx{\MXa+0.06+0.44+0.08}
\node[md block, font=\fontsize{7}{8}\selectfont\bfseries] at (\qx,1.27)
  {Profile A, B or\\no preference?};
% ---- the loop's two edges
\pgfmathsetmacro\xl{\MXa+0.10}\pgfmathsetmacro\xr{\MXb-0.10}
\draw[md arrow=icInk] (\xr,1.695) -- (\xr,1.405);
\node[md block, anchor=east, align=right] at (\xr-0.05,1.55) {Most informative\\pair};
\draw[md arrow=icInk] (\xl,1.405) -- (\xl,1.695);
\node[md block] at (\xl+0.05,1.55) {Belief\\update};
% ---- right: identified assistance -> validation
\draw[md arrow] (\MXb+0.005,2.01) -- (\RXa-0.005,2.01);
\draw[md box] (\RXa,1.70) rectangle (\RXb,\TYb);
\node[md block] at (\RXa+0.06,2.2024)
  {\textcolor{icBlue}{Identified}\\\textcolor{icBlue}{assistance}\\after $k$\\comparisons};
\draw[md arrow] (\rc,1.695) -- (\rc,1.615);
\draw[md box] (\RXa,\TYa) rectangle (\RXb,1.61);
\node[md text, font=\scriptsize\bfseries] at (\RXa+0.06,1.515) {Validation:};
\node[md block] at (\RXa+0.06,1.30) {identified vs.\\generic\\assistance};
% ---- bottom: this paper's simulated users
% the worked answer: the plots' bottom and height, the pair's left edges and width, the answer
% model's left edge and width, the row of labels below them, how far the box grows downward (in)
\def\WEy{0.07}\def\WEh{0.34}\def\WEa{1.07}\def\WEb{1.68}\def\WEw{0.52}
\def\SGx{2.38}\def\SGw{0.54}\def\SGy{\WEy}\def\SGh{\WEh}
\def\WEl{-0.025}\def\WEd{0.12}
\draw[draw=icMuted, line width=\mds*0.8pt, dash pattern=on \mds*2.4pt off \mds*1.6pt,
      rounded corners=0.06*\mdu] (\LXa,0.03-\WEd) rectangle (\RXb,0.70);
\node[md head] at (\LXa+0.07,0.615) {This paper: simulated users};
% who the users are; how they answer: the pair's distances, then the probability they give
\node[md block, anchor=north west] at (\LXa+0.07,0.535)
  {\textbf{Preferred profile}\\\textit{Shape:} walker not\\\phantom{\textit{Shape:} }in the prior\\\textit{Peaks:} study's range};
\node[md block, anchor=north west] at (\WEa,0.535) {\textbf{Distance} $d$ (Nm)};
\node[md block, anchor=north west] at (2.12,0.535) {\textbf{Answer model:} $P$(choose A)};
\node[md text, anchor=base] at (\WEa+\WEw/2,\WEl) {$d(\mathrm{A})\!=\!\icv{fj}{0}{d1_txt}$};
\node[md text, anchor=base] at (\WEb+\WEw/2,\WEl) {$d(\mathrm{B})\!=\!\icv{fj}{0}{d2_txt}$};
\node[md text, anchor=base] at (\SGx+\SGw/2,\WEl) {$d(\mathrm{B}) - d(\mathrm{A})$};
\draw[md arrow] (\WEb+\WEw+0.03,\WEy+\WEh/2) -- (\SGx-0.05,\WEy+\WEh/2);
\draw[md arrow] (\MXa+0.28,0.705) -- (\MXa+0.28,0.795);
% ---- the walkers: three without, the user with the hip exoskeleton
\icfwalker{(\LXa+0.20)*\mdu}{2.315*\mdu}{0.12*\mdu}{5}
\icfwalker{(\LXa+0.38)*\mdu}{2.315*\mdu}{0.12*\mdu}{5}
\icfwalker{(\LXa+0.56)*\mdu}{2.315*\mdu}{0.12*\mdu}{5}
\icfwalker{(\MXa+0.06+0.22-0.02)*\mdu}{(\TYa+0.05+0.02)*\mdu}{0.46*\mdu}{11}
% ---- the small axes
% the prior's shapes at unit peaks: limits of their own
\pgfplotsset{md extra/.style={ymin=\icv{fm}{0}{prior_ymin}, ymax=\icv{fm}{0}{prior_ymax}}}
\mdmini{\LXa+0.15}{1.40}{0.54}{0.29}{%
  \mdpriorfan\addplot[md generic, md line] table[x=phase, y=generic_shape] {\mdprior};}
\pgfplotsset{md extra/.style={}}
\mdmini{\pxa}{\SNy}{0.46}{\SNh}{%
  \mdfan{\mdfanA}{icBlueLight}{0.55}\mdestimate{mean_k0}{1.0}\mdpreferred{0.8}}
\mdmini{\pxb}{\SNy}{0.46}{\SNh}{%
  \mdfan{\mdfanB}{icBlueLight}{0.55}\mdestimate{mean_k6}{1.0}\mdpreferred{0.8}}
\mdmini{\pxc}{\SNy}{0.46}{\SNh}{%
  \mdfan{\mdfanC}{icBlueLight}{0.55}\mdestimate{mean_k24}{1.0}\mdpreferred{0.8}}
% the pair: each label inside its plot's top-left corner, clear of the other plot
\mdmini{\qx}{0.87}{0.46}{0.27}{%
  \addplot[md line, icPairA, line width=\mds*1.2pt] table[x=phase, y=pair1] {\mdpair};
  \node[inner sep=0pt, anchor=north west, font=\scriptsize\bfseries, icInk, xshift=\mds*2.5pt,
        yshift=\mds*1.5pt] at (rel axis cs:0,1) {A};}
\mdmini{\qx+0.60}{0.87}{0.46}{0.27}{%
  \addplot[md line, icPairB, line width=\mds*1.2pt] table[x=phase, y=pair2] {\mdpair};
  \node[inner sep=0pt, anchor=north west, font=\scriptsize\bfseries, icInk, xshift=\mds*2.5pt,
        yshift=\mds*1.5pt] at (rel axis cs:0,1) {B};}
% the same pair with the preferred profile: limits of its own, so that the gaps fill the plots
\pgfplotsset{md extra/.style={ymin=\icv{fj}{0}{ymin}, ymax=\icv{fj}{0}{ymax}}}
\mdmini{\WEa}{\WEy}{\WEw}{\WEh}{%
  \mdgap{pair1}{icPairA}
  \node[inner sep=0pt, anchor=north west, font=\scriptsize\bfseries, icInk, xshift=\mds*2.5pt,
        yshift=\mds*1.5pt] at (rel axis cs:0,1) {A};}
\mdmini{\WEb}{\WEy}{\WEw}{\WEh}{%
  \mdgap{pair2}{icPairB}
  \node[inner sep=0pt, anchor=north west, font=\scriptsize\bfseries, icInk, xshift=\mds*2.5pt,
        yshift=\mds*1.5pt] at (rel axis cs:0,1) {B};}
\pgfplotsset{md extra/.style={}}
\mdmini{\RXa+0.06}{1.75}{\RXb-\RXa-0.12}{0.21}{%
  \mdpreferred{0.9}\mdestimate{mean_k24}{1.2}
  \node[inner sep=0pt, anchor=south east, icMuted]
    at (rel axis cs:0.99,0.02) {$k = \icv{fm}{0}{k_id}$};}
\mdmini{\RXa+0.06}{\TYa+0.05}{\RXb-\RXa-0.12}{0.27}{%
  \mdgeneric\mdestimate{mean_k24}{1.2}}
% ---- the answer model: lapse floor and ceiling, ticks on the right
\pgfmathsetmacro\mdlo{\icv{fm}{0}{lapse}/2}\pgfmathsetmacro\mdhi{1-\icv{fm}{0}{lapse}/2}
\begin{axis}[at={(\SGx*\mdu,\SGy*\mdu)}, anchor=south west, scale only axis,
  width=\SGw*\mdu, height=\SGh*\mdu, clip=false,
  xmin=\icv{fm}{0}{sig_xmin}, xmax=\icv{fm}{0}{sig_xmax},
  ymin=\icv{fm}{0}{sig_ymin}, ymax=\icv{fm}{0}{sig_ymax},
  axis x line*=bottom, axis y line*=right, axis line style={-, icMuted, line width=\mds*0.6pt},
  tick align=outside, major tick length=\mds*2pt, tick style={icMuted, line width=\mds*0.6pt},
  xtick={0}, xticklabels={}, ytick={\mdlo,0.5,\mdhi},
  yticklabels={$\lambda$/2, 0.5, 1\,$-$\,$\lambda$/2}, yticklabel pos=right,
  yticklabel style={icInk, font=\scriptsize, inner sep=0pt, xshift=\mds*1.5pt},
]
\draw[icContext, line width=\mds*0.6pt, dash pattern=on \mds*0.6pt off \mds*0.9pt]
  (axis cs:\icv{fm}{0}{sig_xmin},\mdlo) -- (axis cs:\icv{fm}{0}{sig_xmax},\mdlo);
\draw[icContext, line width=\mds*0.6pt, dash pattern=on \mds*0.6pt off \mds*0.9pt]
  (axis cs:\icv{fm}{0}{sig_xmin},\mdhi) -- (axis cs:\icv{fm}{0}{sig_xmax},\mdhi);
\draw[icRule, line width=\mds*0.5pt]
  (axis cs:0,\icv{fm}{0}{sig_ymin}) -- (axis cs:0,\icv{fm}{0}{sig_ymax});
\addplot[icInk, line width=\mds*1.2pt, mark=none, line join=round] table[x=delta, y=p]
  {\mdsigmoid};
% the pair's reading: its distance difference, the probability it gives
\draw[icMuted, line width=\mds*0.6pt, dash pattern=on \mds*0.9pt off \mds*0.9pt]
  (axis cs:\icv{fj}{0}{delta},\icv{fm}{0}{sig_ymin}) -- (axis cs:\icv{fj}{0}{delta},\icv{fj}{0}{p})
  -- (axis cs:\icv{fm}{0}{sig_xmin},\icv{fj}{0}{p});
\fill[icInk] (axis cs:\icv{fj}{0}{delta},\icv{fj}{0}{p}) circle[radius=\mds*1.3pt];
\node[inner sep=0pt, anchor=north west, font=\scriptsize\bfseries, icInk, yshift=-\mds*1.2pt]
  at (axis cs:\icv{fm}{0}{sig_xmin},\icv{fj}{0}{p}) {\icv{fj}{0}{p_txt}};
\end{axis}
\pgfresetboundingbox
\useasboundingbox (0,-\WEd) rectangle (3.5,2.6);
\end{tikzpicture}%
}
\end{document}
